\documentclass[11pt]{amsart}
\usepackage[margin=1.1in]{geometry}
\usepackage{amsmath,amssymb,amsthm}
\usepackage{enumitem}
\usepackage{booktabs}
\usepackage[dvipsnames]{xcolor}
\usepackage[colorlinks=true,linkcolor=MidnightBlue,citecolor=MidnightBlue,urlcolor=MidnightBlue]{hyperref}

\newtheorem{theorem}{Theorem}[section]
\newtheorem{proposition}[theorem]{Proposition}
\newtheorem{lemma}[theorem]{Lemma}
\newtheorem{corollary}[theorem]{Corollary}
\theoremstyle{definition}
\newtheorem{definition}[theorem]{Definition}
\newtheorem{example}[theorem]{Example}
\newtheorem{remark}[theorem]{Remark}
\newtheorem{question}[theorem]{Question}

% ---- macros -----------------------------------------------------------------
\newcommand{\R}{\mathcal{R}}
\newcommand{\Lr}{\mathcal{L}}
\newcommand{\D}{\mathcal{D}}
\newcommand{\Hh}{\mathcal{H}}
\renewcommand{\P}{\mathcal{P}}
\newcommand{\supp}{\operatorname{supp}}
\newcommand{\T}{\mathsf{T}}
\newcommand{\N}{\mathsf{N}}
\newcommand{\Sym}{\operatorname{Sym}}
\newcommand{\rk}{\operatorname{rank}}
\newcommand{\Z}{\mathbb{Z}}
\newcommand{\Id}{\operatorname{Id}}
\newcommand{\Comp}{\operatorname{Comp}}
\newcommand{\mt}{\bowtie}      % the matched (friendly) relation; replace by \mathcal F if desired
\newcommand{\dn}{\mathord\downarrow}
\newcommand{\todo}[1]{\textcolor{red}{[#1]}}

% ---- condition labels (research label -> paper label; see labels.md) ---------
\newcommand{\Qone}{\textup{(Q1)}}     % projection condition
\newcommand{\Tri}{\textup{(Tri)}}
\newcommand{\Wire}{\textup{(S)}}
\newcommand{\Rcond}{\textup{(R)}}
\newcommand{\Rp}{\textup{(R$'$)}}
\newcommand{\Rpp}{\textup{(R$''$)}}
\newcommand{\Vcond}{\textup{(V)}}
\newcommand{\Pone}{\textup{(P$_1$)}}
\newcommand{\Ptwo}{\textup{(P$_2$)}}
\newcommand{\Qcond}{\textup{(Q)}}
\newcommand{\Eone}{\textup{(E$_1$)}}
\newcommand{\Esix}{\textup{(E$_6$)}}
\newcommand{\Kzero}{\textup{(K0)}}
\newcommand{\Kfour}{\textup{(K4)}}
\newcommand{\Wn}[1][n]{\textup{(W$_{#1}$)}}

\title[Embedding projection algebras in partition monoids]{Which projection algebras embed in a partition monoid?\\ Wiring obstructions and the first non-embeddable examples}
\author{M.~Kinyon}
\date{Draft of \today}

\begin{document}
\maketitle

\begin{abstract}
A regular $*$-semigroup $S$ embeds in a partition monoid $\P_X$ only if its projection algebra $\P(S)$ satisfies the projection condition: projections with a common lower bound commute. We show that this condition is far from sufficient, and that the further obstructions are of a single kind, which we call wiring: a matched path $a\mt c\mt b$ of projections (meaning $aca=a$, $cac=c$, and likewise for $b,c$) which the geometry of $\P_X$ forces to close, $a\mt b$, as soon as auxiliary projections commute with, lie below, or ``see alike'' suitable vertices of the path. We prove eight such quasi-identities valid in every $\P_X$, and an infinite family, the chain conditions, which shows that the quasivariety generated by the projection algebras of partition monoids is not finitely based. Each condition is violated by an explicit finite projection algebra satisfying all the others, and the eight conditions together with the projection condition and the triangle condition characterise embeddability for projection algebras of order at most $8$: there are exactly nine minimal non-embeddable projection algebras of order at most $8$, and we describe them. On the other side we show that the local structure imposes no restriction at all, every finite completely simple regular $*$-semigroup being embeddable, and that for graph semigroups the projection condition and a single wiring condition are also sufficient.
\end{abstract}

\tableofcontents

%==============================================================================
\section{Introduction}\label{sec:intro}

\todo{Phase 4: the introduction is written last. The present text fixes the statements of the main results and the vocabulary of the paper.}

Let $X$ be a set. The partition monoid $\P_X$ consists of the set partitions of $X\sqcup X'$, where $X'$ is a disjoint copy of $X$, under the usual product by concatenation; with the involution $\alpha\mapsto\alpha^*$ that interchanges $X$ and $X'$ it is a regular $*$-semigroup, and it contains the symmetric group, the symmetric inverse monoid, the Brauer and Temperley--Lieb monoids and many others as $*$-subsemigroups. The question we study is which regular $*$-semigroups embed in some $\P_X$ by a $*$-homomorphism, and more specifically which \emph{projection algebras} do. Recall that the projections of a regular $*$-semigroup are the elements $p=p^2=p^*$, and that they form a \emph{projection algebra} under the binary operation $\theta_p(q)=pqp$; by a theorem of Imaoka every abstract projection algebra, in the axiomatic sense of Jones, is the projection algebra of some regular $*$-semigroup, and the projection algebra determines much of the structure of $S$ \todo{refs: Imaoka, Jones, East--Gray--Azeef Muhammed--Ru\v skuc}. A $*$-embedding $S\to\P_X$ restricts to an embedding of projection algebras $\P(S)\to\P(\P_X)$, and conversely an embedding $\P(S)\to\P(\P_X)$ extends to $S$ whenever $S$ is generated by its projections \todo{precise statement and reference}. So the question for projection algebras is the heart of the matter, and it is the one we address.

One necessary condition has been known: in $\P_X$, two projections with a common lower bound in the natural order commute, i.e.\ the quasi-identity
\[
\Qone\qquad r=prp,\ r=qrq\ \Longrightarrow\ pq=qp
\]
holds in $\P(\P_X)$ \todo{ref}. Our results are the following.

\emph{The projection condition is far from sufficient, and the further obstructions are wiring conditions.} All the conditions we find have the same shape. Call projections $p,q$ \emph{matched}, $p\mt q$, if $pqp=p$ and $qpq=q$; this is the \emph{friendly} relation of \todo{EGMR}, and in $\P_X$ it means that $p$ and $q$ have the same rank and their transversal blocks are paired off by the product. A matched path $a\mt c\mt b$ need not close ($a\mt b$ need not hold), but in $\P_X$ it is forced to close by any of a number of side conditions on auxiliary projections: that an auxiliary commutes with two vertices and sees the third as it sees one of them, that two vertices see an auxiliary alike, that an auxiliary lies above a vertex, and so on. Each such side condition is a quasi-identity in the language of projection algebras, and we prove eight of them in every $\P_X$ (Theorem~\ref{thm:main-wiring}), by arguments that analyse the blocks of the projections involved. The simplest, \Wire, has four variables and explains the smallest non-embeddable example, a projection algebra of order $6$; the others have up to seven variables.

\emph{The conditions are not finitely many.} The wiring conditions extend to an infinite family, the \emph{chain conditions} \Wn, in which the auxiliary is replaced by a chain $c\sim x_1\sim\dots\sim x_n\sim b$ of pairwise commuting projections each of which identifies $a$ with $c$; we prove all of them in every $\P_X$ by a counting invariant (Theorem~\ref{thm:main-chain}), and we exhibit for each $n$ a projection algebra of order $2n+5$ that satisfies \Wn[n-1] and all the other conditions of this paper but violates \Wn. Consequently the quasivariety generated by the projection algebras of partition monoids is not finitely based, and the class of finite embeddable projection algebras is not the class of finite models of any finite set of quasi-identities (Theorem~\ref{thm:main-nfb}).

\emph{The small examples are classified.} Up to order $8$ the obstructions are completely understood: a projection algebra of order at most $8$ embeds in some $\P(\P_X)$ if and only if it satisfies \Qone\ together with the wiring conditions of Theorem~\ref{thm:main-wiring}, and there are exactly nine minimal non-embeddable projection algebras of order at most $8$, one of order $6$, one of order $7$ and seven of order $8$, each the witness of one condition (Theorem~\ref{thm:main-census}). The computations behind this (an enumeration of projection algebras and of regular $*$-semigroups up to order $9$, and embedding searches into $\P(\P_m)$ for $m\le8$) are described in Section~\ref{sec:comp}.

\emph{The local structure is no obstruction, and for graph semigroups the conditions are sufficient.} Every finite completely simple regular $*$-semigroup embeds in a finite partition monoid (Theorem~\ref{thm:main-cs}), so no condition on the principal factors can be necessary; and for the graph semigroups $Q(G,\bot)$ of \todo{ref}, and for a layered generalisation of them, \Qone\ and \Wire\ are also sufficient (Theorem~\ref{thm:main-graph}).

\emph{The digraph conditions.} Besides the wiring conditions there is a family of conditions of a different shape, indexed by digraphs, whose simplest member is the \emph{triangle condition} \Tri: three pairwise commuting projections with common product $w$, three with common product $z$, matched crosswise, force $w=z$. In every projection algebra the triangle condition implies all the others of its family; it holds in every $\P_X$, and it fails in a $61$-element quotient of $\P_3$ that satisfies \Qone\ (Theorem~\ref{thm:main-digraph}). \todo{Decide in Phase 2 how much of Section~\ref{sec:digraph} to keep; the infinite case is compressed to a remark.}

The methods are elementary but the proofs are long, because the only tool is the block calculus of $\P_X$: a projection of $\P_X$ is a partition of $X$ with some of its classes marked, a product of projections is described by the join of their partitions, and a quasi-identity is proved by following chains of blocks. Two general devices organise these arguments: a \emph{counting invariant} attached to pairs of projections, which is constant along chains of commuting identifiers (Section~\ref{sec:chain}), and a \emph{finite descent} principle (Lemma~\ref{lem:descent}), which replaces the rank of $\P_X$ by the size of a principal down-set in an arbitrary finite projection algebra and lets several arguments be finished algebraically.

\medskip
\noindent\emph{Main theorems.} \todo{Numbering and cross-references to be fixed once the sections are in place; for now the theorems are stated here and proved in the sections indicated.}

\begin{theorem}[local structure; Section~\ref{sec:local}]\label{thm:main-cs}
Every finite completely simple regular $*$-semigroup embeds in a finite partition monoid. $\P_4$ contains non-orthodox completely simple $*$-subsemigroups.
\end{theorem}

\begin{theorem}[digraph conditions; Section~\ref{sec:digraph}]\label{thm:main-digraph}
For every admissible digraph $E$ the condition $(\mathrm D_E)$ holds in every $\P_X$. In every projection algebra the triangle condition implies $(\mathrm D_E)$ for every $E$. The triangle condition fails in a $61$-element regular $*$-semigroup satisfying \Qone, a quotient of $\P_3$, and in no projection algebra of order at most $10$ satisfying \Qone.
\end{theorem}

\begin{theorem}[wiring conditions; Section~\ref{sec:wiring}]\label{thm:main-wiring}
The conditions \Wire, \Rcond, \Vcond, \Rp, \Pone, \Qcond, \Kfour\ and \Kzero\ of Section~\ref{sec:wiring} hold in every $\P_X$. Each is violated by an explicit projection algebra of order at most $8$ which satisfies \Qone\ and all the other conditions of this paper.
\end{theorem}

\begin{theorem}[chain conditions; Section~\ref{sec:chain}]\label{thm:main-chain}
For every $n\ge1$ the chain condition \Wn\ holds in every $\P_X$. For every $n$ there is a projection algebra $\mathfrak P^{\mathrm{ch}}_n$ of order $2n+5$ which satisfies \Wn[n-1], \Qone, \Tri\ and the eight wiring conditions, and violates \Wn.
\end{theorem}

\begin{theorem}[no finite basis]\label{thm:main-nfb}
The quasivariety generated by the projection algebras $\P(\P_X)$, $X$ finite, is not finitely based; the class of finite projection algebras embeddable in some $\P(\P_X)$ is not the class of finite models of any finite set of quasi-identities.
\end{theorem}

\begin{theorem}[orders at most $8$; Sections~\ref{sec:wiring} and \ref{sec:comp}]\label{thm:main-census}
A projection algebra of order at most $8$ embeds in some $\P(\P_X)$ if and only if it satisfies \Qone, \Tri\ and the eight wiring conditions. The minimal non-embeddable projection algebras of order at most $8$ are exactly the nine algebras $\mathfrak P$ (order $6$), $\mathfrak P_7$ (order $7$) and $\mathfrak P_8^1,\dots,\mathfrak P_8^7$ (order $8$) of Section~\ref{sec:wiring}.
\end{theorem}

\begin{theorem}[graph semigroups; Section~\ref{sec:graph}]\label{thm:main-graph}
A graph semigroup $Q(G,\bot)$, and more generally a layered graph semigroup, embeds in a partition monoid if and only if it satisfies \Qone\ and \Wire.
\end{theorem}

\medskip
\noindent\emph{Evidence levels.} Throughout we distinguish three kinds of statement: a quasi-identity is \emph{proved} when it is shown to hold in every $\P_X$; it is \emph{verified for $m\le k$} when a SAT solver has shown it to hold in $\P(\P_m)$ for all $m\le k$ (with variables allowed to coincide); and a finite algebra is \emph{non-embeddable for $m\le k$} when a search has found no embedding into $\P(\P_m)$ for $m\le k$. Only the first is a theorem; the others are reported as such, and the main results use only proved statements.

%==============================================================================
\section{Preliminaries}\label{sec:prelim}

\subsection{Regular $*$-semigroups and projection algebras}\label{subsec:rss}

A \emph{regular $*$-semigroup} is a semigroup $S$ with a unary operation $x\mapsto x^*$ satisfying $x^{**}=x$, $(xy)^*=y^*x^*$ and $xx^*x=x$. Its \emph{projections} are the elements $p=p^2=p^*$; we write $\P(S)$ for the set of projections. For projections $p,q$ the element $\theta_p(q)=pqp$ is again a projection, and $\P(S)$ with the binary operation $(p,q)\mapsto\theta_p(q)$ is the \emph{projection algebra} of $S$. In the notation $p\cdot q:=\theta_q(p)$ (the value lies below the second argument) the identities satisfied by all projection algebras of regular $*$-semigroups are
\begin{align}
&p\cdot p=p, \qquad (q\cdot p)\cdot p=q\cdot p, \qquad (p\cdot q)\cdot p=q\cdot p,\label{eq:PA123}\\
&((r\cdot p)\cdot q)\cdot p=r\cdot(q\cdot p), \qquad (((r\cdot p)\cdot q)\cdot p)\cdot q=(r\cdot p)\cdot q,\label{eq:PA45}
\end{align}
and conversely an algebra $(P,\cdot)$ satisfying \eqref{eq:PA123}--\eqref{eq:PA45} is the projection algebra of some regular $*$-semigroup \todo{Imaoka 1983; Jones 2012; cite the exact axiom system used}. In $\theta$-notation the axioms read $\theta_p(p)=p$, $\theta_p\theta_p=\theta_p$, $\theta_p(\theta_q(p))=\theta_p(q)$, $\theta_p\theta_q\theta_p=\theta_{\theta_p(q)}$ and $\theta_q\theta_p\theta_q\theta_p=\theta_q\theta_p$. We call an abstract algebra satisfying them a projection algebra; everything below about projection algebras uses only these axioms, except where finiteness is invoked explicitly.

The \emph{natural order} on projections is $q\le p$ iff $q=pq=qp$, equivalently $q=\theta_p(q)$. We write $\dn p=\{q:q\le p\}$. Two projections \emph{commute}, $p\sim q$, if $pq=qp$; then $pq$ is a projection, $pq\le p,q$, and $pq=\theta_p(q)=\theta_q(p)$. In the projection algebra alone commutation is expressed by $\theta_p(q)=\theta_q(p)$:

\begin{lemma}\label{lem:comm}
For projections $p,q$ of a regular $*$-semigroup, $pq=qp$ iff $\theta_p(q)=\theta_q(p)$.
\end{lemma}

\begin{proof}
If $pq=qp$ then $pqp=pq=qpq$. Conversely let $w=pqp=qpq$. Then $(pq)(pq)^*=pqqp=w$, so by regularity $pq=(pq)(pq)^*(pq)=w\,pq=pqp\,pq=(pq)^2=p(qpq)=pw=p\,pqp=w$; symmetrically $qp=w$.
\end{proof}

We note in passing that the product of two projections is always idempotent, since $pq=(pq)(pq)^*(pq)=pqqppq=pqpq$; in general it is not a projection.

\begin{definition}[matched pairs]\label{def:matched}
Projections $p,q$ are \emph{matched}, $p\mt q$, if $\theta_p(q)=p$ and $\theta_q(p)=q$.
\end{definition}

This is the \emph{friendly} relation $p\,\mathcal F\,q$ of \todo{EGMR}; we keep a separate symbol because in $\P_X$ it comes with a bijection of transversal blocks (Lemma~\ref{lem:PX}(c)) and the word ``matched'' refers to that bijection. The relation is reflexive and symmetric, not transitive, and matched pairs can be manufactured at will:

\begin{lemma}\label{lem:friends}
For all projections $p,q$: $\theta_p(q)\mt\theta_q(p)$, and $\theta_{\theta_p(q)}(q)=\theta_p(q)$.
\end{lemma}

\begin{proof}
Both statements follow from the axioms in $\theta$-notation. Using $\theta_p\theta_q\theta_p=\theta_{\theta_p(q)}$ and then $\theta_p(\theta_q(p))=\theta_p(q)$ twice,
\[
\theta_{\theta_p(q)}(\theta_q(p))=\theta_p\theta_q\theta_p(\theta_q(p))=\theta_p\theta_q(\theta_p(q))=\theta_p(\theta_q(p))=\theta_p(q),
\]
and symmetrically $\theta_{\theta_q(p)}(\theta_p(q))=\theta_q(p)$; so $\theta_p(q)\mt\theta_q(p)$. Likewise $\theta_{\theta_p(q)}(q)=\theta_p\theta_q\theta_p(q)=\theta_p(\theta_q(p))=\theta_p(q)$.
\end{proof}

\subsection{The projection condition}\label{subsec:Q1}

In $\P_X$, as we recall in Section~\ref{subsec:PX}, two projections with a common lower bound commute. The corresponding quasi-identity
\begin{equation}\label{eq:Q1}\tag{Q1}
\theta_p(r)=r=\theta_q(r)\ \Longrightarrow\ \theta_p(q)=\theta_q(p)
\end{equation}
is therefore necessary for embeddability. It has several equivalent forms.

\begin{proposition}\label{prop:Q1}
For a regular $*$-semigroup $S$ the following are equivalent.
\begin{enumerate}[label=\textup{(\arabic*)}]
\item $\P(S)$ satisfies \eqref{eq:Q1}.
\item Every idempotent of $S$ lying above a projection in the natural order is a projection.
\item For every projection $p$, the set $S_p=\{x\in S: px=xp=p\}$ is an inverse semigroup.
\item Two projections commute if and only if they have a common lower bound.
\end{enumerate}
\end{proposition}

\begin{proof}
$S_p$ is closed under multiplication and under $*$ (if $px=xp=p$ then $px^*=(xp)^*=p$), so it is a regular $*$-semigroup with zero $p$; it is inverse iff all its idempotents are projections, and its idempotents are exactly the idempotents $e\ge p$. This gives (2)$\Leftrightarrow$(3).

(2)$\Rightarrow$(1). Let $r\le p,q$. The element $e=pq$ is idempotent, and $er=pqr=pr=r=re$, so $e$ lies above the projection $r$; by (2) it is a projection, and $pq=e=e^*=qp$.

(1)$\Rightarrow$(2). Let $e=e^2$ and let $r$ be a projection with $er=re=r$. Then $r=r^*=(re)^*=e^*r$ and likewise $r=re^*$, so the projections $f=ee^*$ and $g=e^*e$ both lie above $r$: $fr=ee^*r=er=r$, $rf=re\,e^*=re^*=r$, and similarly for $g$. By (1), $fg=gf$, that is, $ee^*e^*e=e^*ee^*$. Multiplying this on the left by $e$ gives $fg=e\,e^*e^*e=(ee^*)(ee^*)=f$, and multiplying it on the right by $e$ gives $fg=f(ge)=(fg)e=(gf)e=e^*e\,ee^*\,e=(e^*e)(e^*e)=g$ (using $ge=e^*ee=g$). Hence $ee^*=e^*e$, and $e=ee^*e=e^*ee=e^*e$ is a projection.

(1)$\Leftrightarrow$(4). If $pq=qp$ then $pq$ is a projection below $p$ and $q$; the converse is \eqref{eq:Q1} with Lemma~\ref{lem:comm}.
\end{proof}

\begin{remark}\label{rem:Sp}
In $\P_X$, an element $x$ lies above a projection $p$ iff $x$ maps every $\T(p)$-block onto itself by a block bijection and permutes the $\N(p)$-blocks, so that $(\P_X)_p\cong\prod_{A\in\T(p)}\mathcal I^*_A\times\mathcal J_p$ with $\mathcal I^*_A$ the dual symmetric inverse monoid on $A$ and $\mathcal J_p$ an inverse monoid of partial block bijections; for the projection $p$ with all blocks singletons and nontransversal, $(\P_X)_p$ is the symmetric inverse monoid $\mathcal I_X$. So \eqref{eq:Q1} says that these ``unitary'' subsemigroups are inverse.
\end{remark}

Two consequences of \eqref{eq:Q1} are used throughout: the projections above a fixed projection pairwise commute; and two distinct matched projections have no common lower bound (they would commute, and then $p=pqp=pq=q$).

\subsection{Partition monoids and the block calculus}\label{subsec:PX}

For a set $X$ let $X'=\{x':x\in X\}$ be a disjoint copy. The partition monoid $\P_X$ consists of the set partitions of $X\sqcup X'$; the product $\alpha\beta$ is obtained by identifying the lower points of $\alpha$ with the upper points of $\beta$, taking connected components, and forgetting the middle layer, and $\alpha^*$ interchanges $X$ and $X'$. A block is \emph{transversal} if it meets both $X$ and $X'$; the \emph{rank} of $\alpha$ is the number of transversal blocks.

A projection $p$ of $\P_X$ is determined by a pair $(\varepsilon_p,\T(p))$, where $\varepsilon_p$ is a partition of $X$ and $\T(p)$ is a set of $\varepsilon_p$-classes: the blocks of $p$ are $A\cup A'$ for $A\in\T(p)$ (the \emph{$\T$-blocks}) and $B$, $B'$ for the remaining classes $B$ (the \emph{$\N$-blocks}, forming $\N(p)$). We put $\supp p=\bigcup\T(p)$, so that $\rk p=|\T(p)|$. We say ``block of $p$'' for a class of $\varepsilon_p$, with the label $\T$ or $\N$ understood. For projections $p,q$ the \emph{join} $\varepsilon_p\vee\varepsilon_q$ is the join of partitions of $X$; its classes are the connected components of the graph on the blocks of $p$ and $q$ in which two blocks are adjacent when they meet. A \emph{chain} in $\varepsilon_p\vee\varepsilon_q$ is a sequence of blocks, alternately of $p$ and of $q$, in which consecutive blocks meet; two blocks lie in a common class iff they are joined by a chain. For a class $K$ we write $U_K$ for the union of the $\T(p)$-blocks contained in $K$ and $L_K$ for the union of the $\T(q)$-blocks in $K$. The following lemma is the whole of our toolkit.

\begin{lemma}\label{lem:PX}
Let $p,q$ be projections of $\P_X$.
\begin{enumerate}[label=\textup{(\alph*)}]
\item The blocks of $\theta_p(q)=pqp$ are the $\N(p)$-blocks and the sets $U_K$ for the classes $K$ of $\varepsilon_p\vee\varepsilon_q$ meeting $\supp p$; $U_K$ is a $\T$-block of $\theta_p(q)$ iff $K$ meets $\supp q$. In particular the blocks of $\theta_p(q)$ not contained in $\supp p$ are exactly the $\N(p)$-blocks; consequently, if $e\le p$ then $\supp e\subseteq\supp p$, every $\T(e)$-block is a union of $\T(p)$-blocks, and every block of $e$ not contained in $\supp p$ is an $\N(p)$-block.
\item Suppose $pq=qp$. Then every class $K$ of $\varepsilon_p\vee\varepsilon_q$ is of one of four types: \textup{(i)} $K=U_K=L_K$; \textup{(ii)} $K=U_K$ is a single $\N(q)$-block; \textup{(iii)} $K=L_K$ is a single $\N(p)$-block; \textup{(iv)} $K$ is a single block of both $p$ and $q$. The $\T$-blocks of $pq$ are the classes of type \textup{(i)}; hence $\supp(pq)=\supp p\cap\supp q$, and every $\N(q)$-block contained in $\supp p$ is a union of $\T(p)$-blocks.
\item Suppose $p\mt q$. Then every class of $\varepsilon_p\vee\varepsilon_q$ meeting $\supp p\cup\supp q$ contains exactly one $\T(p)$-block and exactly one $\T(q)$-block. In particular $\rk p=\rk q$, no $\T(p)$-block is an $\N(q)$-block, and the classes of $\varepsilon_p\vee\varepsilon_q$ induce a bijection $\T(p)\to\T(q)$. Conversely, if every class meeting $\supp p\cup\supp q$ contains exactly one $\T(p)$-block and one $\T(q)$-block then $p\mt q$.
\end{enumerate}
\end{lemma}

\begin{proof}
(a) In the product $pq$ the middle layer is partitioned by $\varepsilon_p\vee\varepsilon_q$, since every block of $p$ has a lower part and every block of $q$ an upper part in the middle; the transversal blocks of $pq$ are $U_K\cup L_K'$ for the classes $K$ with $U_K\ne\emptyset\ne L_K$, the upper nontransversal blocks are the $\N(p)$-blocks and the sets $U_K$ with $L_K=\emptyset$. Now $\theta_p(q)=(pq)(pq)^*$ is the projection whose $\T$-blocks are the upper parts of the transversal blocks of $pq$ and whose $\N$-blocks are the upper nontransversal blocks of $pq$. For the last sentence use $e=\theta_p(e)$.

(b) The projection $pq$ is symmetric, so $U_K\cup L_K'$ is a block only if $U_K=L_K$ (when both are nonempty), and the family of upper nontransversal blocks $\N(p)\cup\{U_K:L_K=\emptyset\ne U_K\}$ equals the family of lower ones $\N(q)\cup\{L_K:U_K=\emptyset\ne L_K\}$. Let $K$ be a class. If $U_K\ne\emptyset\ne L_K$, an $\N(p)$-block inside $K$ would be an $\N$-block of $pq$, hence (as $U_K\ne\emptyset$) an $\N(q)$-block; but a set which is a block of both $p$ and $q$ is a class by itself; so $K$ contains no $\N(p)$-block and dually no $\N(q)$-block, and $K=U_K=L_K$, type (i). If $U_K\ne\emptyset=L_K$ then $U_K$ is an $\N(q)$-block; $K$ is a union of $q$-blocks containing the $q$-block $U_K$ and $p$-blocks; as before the remaining $p$-blocks in $K$ would be common blocks, so $K=U_K$, type (ii). Type (iii) is dual, and if $U_K=L_K=\emptyset$ then $K$ consists of $\N(p)$- and $\N(q)$-blocks, which must coincide as above, type (iv). Finally an $\N(q)$-block $B\subseteq\supp p$ lies in a class $K$ which is not of type (i) or (iii) (these consist of $\T(q)$-blocks only) nor of type (iv) ($B$ meets $\supp p$), so $K=B$ is of type (ii).

(c) $pqp=p$ says that the upper projection of $pq$ is $p$: its $\T$-blocks $\{U_K:U_K\ne\emptyset\ne L_K\}$ are exactly the $\T(p)$-blocks, and its $\N$-blocks $\N(p)\cup\{U_K:U_K\ne\emptyset=L_K\}$ are exactly the $\N(p)$-blocks. So every class meeting $\supp p$ contains exactly one $\T(p)$-block and at least one $\T(q)$-block; by symmetry, exactly one of each. A $\T(p)$-block which is an $\N(q)$-block would be a class by itself, containing no $\T(q)$-block. The converse is read off from (a).
\end{proof}

We shall use (a) constantly in the following form: for projections $p,q$, \emph{$p$ sees $q$ as} $\theta_p(q)$, and $\theta_p(q)$ is determined by the partition of $\supp p$ into the sets $U_K$ together with the information which of them are \emph{bright} (meet $\supp q$) and which are \emph{dark}. Thus ``$p$ sees $q$ and $r$ alike'', $\theta_p(q)=\theta_p(r)$, means that $\varepsilon_p\vee\varepsilon_q$ and $\varepsilon_p\vee\varepsilon_r$ induce the same partition of $\supp p$ with the same bright parts; we then say that $p$ \emph{identifies} $q$ with $r$ and write $p\in\Id(q,r)$. Two orders of magnitude of this statement will matter: $\theta_p(q)=p$ (every class meeting $\supp p$ contains one $\T(p)$-block and is bright: $p$ is \emph{half-matched} with $q$) and $\theta_p(q)=q$ ($q\le p$).

\begin{lemma}\label{lem:order}
Let $p,q$ be projections of $\P_X$.
\begin{enumerate}[label=\textup{(\alph*)}]
\item $q\le p$ iff every block of $p$ lies in a block of $q$ and every $\N(p)$-block is a block of $q$. In that case every block of $q$ is an $\N(p)$-block or a union of $\T(p)$-blocks, and every $\N(p)$-block is an $\N(q)$-block.
\item $\theta_p(q)=p$ iff every class of $\varepsilon_p\vee\varepsilon_q$ meeting $\supp p$ contains exactly one $\T(p)$-block and meets $\supp q$. In particular every block of $q$ then meets at most one $\T(p)$-block.
\end{enumerate}
\end{lemma}

\begin{proof}
(a) If $q=\theta_p(q)$, Lemma~\ref{lem:PX}(a) describes the blocks of $q$ as stated. Conversely, if the two conditions hold then $\varepsilon_p\vee\varepsilon_q=\varepsilon_q$, every class of $\varepsilon_q$ meeting $\supp p$ is a union of $\T(p)$-blocks, and it is a $\T(q)$-block iff it meets $\supp q$; so $\theta_p(q)=q$ by Lemma~\ref{lem:PX}(a). (b) is Lemma~\ref{lem:PX}(a) read for $\theta_p(q)=p$.
\end{proof}

For a matched pair $a\mt c$ we write $A_1,\dots,A_m$ for the $\T(a)$-blocks, $C_i$ for the $\T(c)$-block in the class of $A_i$, and $\Gamma_i$ for that class of $\varepsilon_a\vee\varepsilon_c$; so $\Gamma_i\cap\supp a=A_i$, $\Gamma_i\cap\supp c=C_i$, and the other classes of $\varepsilon_a\vee\varepsilon_c$ meet neither support. For a matched path $a\mt c\mt b$ we index the $\T(b)$-blocks through $c$: $B_i$ is the $\T(b)$-block in the class $\Delta_i$ of $\varepsilon_b\vee\varepsilon_c$ containing $C_i$. We also put
\[
u=\theta_a(b),\qquad v=\theta_b(a),
\]
so that $u\le a$, $v\le b$, $u\mt v$ (Lemma~\ref{lem:friends}), and $a\mt b$ iff $u=a$ and $v=b$. By Lemma~\ref{lem:PX}(a) the $\T$-blocks of $u$ are the sets $D\cap\supp a$ for the classes $D$ of $\varepsilon_a\vee\varepsilon_b$ meeting both supports, and $a\mt b$ iff every class of $\varepsilon_a\vee\varepsilon_b$ meeting $\supp a\cup\supp b$ contains exactly one $A$ and one $B$.

\subsection{Finite descent}\label{subsec:descent}

In $\P_X$ many arguments end with a rank count: a projection below a projection of the same rank equals it. The following lemma provides the same conclusion in every finite projection algebra, with the size of a principal down-set in place of the rank.

\begin{lemma}[finite descent]\label{lem:descent}
Let $P$ be a projection algebra.
\begin{enumerate}[label=\textup{(\alph*)}]
\item If $p\mt q$ then $\theta_q$ maps $\dn p$ into $\dn q$, $\theta_p$ maps $\dn q$ into $\dn p$, and the two maps are mutually inverse bijections.
\item If $P$ is finite, $q\le p$, and $p,q$ lie in one component of the graph of $\mt$, then $q=p$.
\item If $P$ is finite, $\theta_p(a)=p$, and $a,p$ lie in one component of the graph of $\mt$, then $a\mt p$.
\item If $P$ is finite and $a\mt c\mt b$, then $a\mt b$ iff $\theta_a(b)=a$ iff $\theta_b(a)=b$.
\end{enumerate}
\end{lemma}

\begin{proof}
(a) $\theta_q(e)\le q$ for every $e$. For $e\le p$ the axiom $\theta_p\theta_q\theta_p=\theta_{\theta_p(q)}$ gives $\theta_p(\theta_q(e))=\theta_p\theta_q\theta_p(e)=\theta_{\theta_p(q)}(e)=\theta_p(e)=e$, and symmetrically. (b) By (a) and induction along a chain $q=x_0\mt x_1\mt\dots\mt x_k=p$, $|\dn q|=|\dn p|$; as $\dn q\subseteq\dn p$, the two sets coincide, so $p\le q$ and $q=p$. (c) Put $a'=\theta_a(p)\le a$. By Lemma~\ref{lem:friends}, $a'\mt\theta_p(a)=p$, so $a'$ lies in the component of $p$, hence of $a$, and (b) gives $a'=a$, i.e.\ $\theta_a(p)=a$. (d) If $u=\theta_a(b)=a$ then $v=\theta_b(a)\le b$ and $v\mt u=a\mt c\mt b$, so $v=b$ by (b); symmetrically.
\end{proof}

In $\P_X$ the down-set $\dn p$ depends only on $\rk p$, so (b) is the rank count; but the lemma holds in every finite projection algebra, and this has consequences for the logical relations among our conditions: a quasi-identity whose proof in $\P_X$ uses only the block calculus and rank counts along matched chains often has a finite-algebraic twin which is not derivable equationally. We shall meet one instance in Section~\ref{subsec:K0}.

%==============================================================================
\section{The local structure is unconstrained}\label{sec:local}

A natural first guess for a condition beyond \eqref{eq:Q1} is a restriction on the ``local'' pieces of $S$, that is, on its principal factors, which are completely $0$-simple regular $*$-semigroups. Some care is needed in saying what such a restriction could be. The zero of a principal factor is not in general an element of $S$, and when it is -- when $S$ contains a completely $0$-simple $*$-subsemigroup together with its zero -- \eqref{eq:Q1} already forces that subsemigroup to be a Brandt semigroup, since two distinct matched projections have no common lower bound (Section~\ref{subsec:Q1}). So the zero pattern of a principal factor is not local information: it is decided by how the $\D$-class sits inside $S$, and it interacts with the rest of $S$, as the triangle condition of Section~\ref{sec:digraph} shows. What remains genuinely local is the nonzero part of a principal factor: a set of projections of a $\D$-class with pairwise nonzero products, together with the sandwich entries and the maximal subgroup. These are exactly the completely simple $*$-subsemigroups of $S$, and a condition on them would be a condition on finite completely simple regular $*$-semigroups. There is none.

Recall that a completely simple regular $*$-semigroup is a Rees matrix semigroup $\mathcal M[G;I,I;P]$ with involution $(i,g,j)^*=(j,g^{-1},i)$, and that the sandwich matrix may be normalised so that $p_{ii}=1$ and $p_{ji}=p_{ij}^{-1}$ for all $i,j\in I$ \todo{ref: Nordahl--Scheiblich}. Its projections are the elements $(i,1,i)$.

\begin{theorem}\label{thm:cs}
Every finite completely simple regular $*$-semigroup embeds in a finite partition monoid. Explicitly, let $S=\mathcal M[G;I,I;P]$ be in normal form, let $G\le\Sym(k)$ act on $[k]=\{1,\dots,k\}$ on the right, and let
\[
X=\{a_{i,s}: i\in I,\ s\in[k]\}\ \sqcup\ \{c_{ij,s}: i\ne j\in I,\ s\in[k]\},\qquad |X|=k|I|^2 .
\]
For $i\in I$ let $\pi_i$ be the partition of $X$ into the singletons $\{a_{i,s}\}$, the triples
\[
T_{i;l,s}=\{c_{il,s},\ a_{l,sp_{il}},\ c_{li,sp_{il}}\}\qquad(l\in I\setminus\{i\},\ s\in[k])
\]
and the singletons $\{c_{lm,s}\}$ with $i\notin\{l,m\}$. Let $\varphi(i,g,j)$ be the partition of $X\sqcup X'$ with transversal blocks $\{a_{i,s},a'_{j,sg}\}$ $(s\in[k])$, upper nontransversal blocks the blocks of $\pi_i$ other than the singletons $\{a_{i,s}\}$, and lower nontransversal blocks the primed copies of the blocks of $\pi_j$ other than the singletons $\{a_{j,s}\}$. Then $\varphi\colon S\to\P_X$ is an injective $*$-homomorphism.
\end{theorem}

\begin{proof}
That $\pi_i$ is a partition of $X$ is checked by locating each point: $a_{l,s}$ with $l\ne i$ lies in $T_{i;l,sp_{li}}$ only, $c_{il,s}$ in $T_{i;l,s}$ only, $c_{li,s}$ in $T_{i;l,sp_{li}}$ only, and the other points are singletons. The blocks of $\pi_i$ other than the $\{a_{i,s}\}$ cover $X\setminus\{a_{i,s}:s\in[k]\}$, so $\varphi(i,g,j)$ is a partition of $X\sqcup X'$, of rank $k$. The description of $\varphi(i,g,j)^*$ is obtained by interchanging $i$ with $j$ and $g$ with $g^{-1}$, so $\varphi$ respects the involution; and $\varphi$ is injective, since the transversal blocks determine $i$, $j$ and the permutation $s\mapsto sg$, hence $g$.

For the product $\varphi(i,g,j)\varphi(l,h,m)$ consider the middle layer, which carries the partition $\pi_j$ (from the lower blocks of the first factor, with the points $a_{j,t}$ attached to the transversal blocks $\{a_{i,s},a'_{j,sg}\}$) and the partition $\pi_l$ (from the upper blocks of the second factor, with the points $a_{l,t}$ attached to the transversal blocks $\{a_{l,t},a'_{m,th}\}$). If $j=l$ the two partitions coincide, so the components of the middle layer are the blocks of $\pi_j$; the point $a_{j,sg}$ is a singleton of $\pi_j$ and is attached to $a_{i,s}$ above and $a'_{m,sgh}$ below, and no other middle block meets a transversal block. The product therefore has transversal blocks $\{a_{i,s},a'_{m,sgh}\}$, upper nontransversal blocks those of the first factor and lower ones those of the second: it is $\varphi(i,gh,m)=\varphi((i,g,j)(j,h,m))$, as $p_{jj}=1$. If $j\ne l$, put $t=sg\,p_{jl}$ and consider the set
\[
M_s=\{a_{j,sg},\ c_{jl,sg},\ c_{lj,t},\ a_{l,t}\}=T_{l;j,t}\cup T_{j;l,sg}\cup\{a_{j,sg}\}\cup\{a_{l,t}\},
\]
where $T_{l;j,t}=\{c_{lj,t},a_{j,tp_{lj}},c_{jl,tp_{lj}}\}=\{c_{lj,t},a_{j,sg},c_{jl,sg}\}$ is a block of $\pi_l$ and $T_{j;l,sg}=\{c_{jl,sg},a_{l,t},c_{lj,t}\}$ is a block of $\pi_j$. Each point of $M_s$ lies, in each of $\pi_j$ and $\pi_l$, in a block contained in $M_s$ (the points $a_{j,sg}$ and $a_{l,t}$ are singletons of $\pi_j$ and $\pi_l$ respectively), so $M_s$ is a component of the middle layer; it contains exactly one transversal end $a_{j,sg}$ of the first factor, attached to $a_{i,s}$, and exactly one transversal start $a_{l,t}$ of the second, attached to $a'_{m,th}$. Hence the product has the transversal blocks $\{a_{i,s},a'_{m,sgp_{jl}h}\}$, which are those of $\varphi(i,gp_{jl}h,m)$. The upper nontransversal blocks of the first factor have no points in the middle layer, so they are blocks of the product; every other upper point lies in a transversal block; so the upper nontransversal blocks of the product are those of $\varphi(i,\cdot,\cdot)$, and dually for the lower ones. Thus $\varphi(i,g,j)\varphi(l,h,m)=\varphi(i,gp_{jl}h,m)=\varphi((i,g,j)(l,h,m))$.
\end{proof}

Since every finite group embeds in some $\Sym(k)$, every finite completely simple regular $*$-semigroup embeds in some $\P_n$. Two remarks on what this rules out.

\begin{remark}\label{rem:orthodox}
Orthodoxy of the completely simple $*$-subsemigroups is not necessary either, although $\P_n$ for small $n$ looks ``flat'': the projections
\begin{align*}
p&: \{0\},\{1\}\ \text{transversal},\ \{2,3\}\ \text{nontransversal};\\
q&: \{1\},\{3\}\ \text{transversal},\ \{0,2\}\ \text{nontransversal};\\
r&: \{0\},\{3\}\ \text{transversal},\ \{1,2\}\ \text{nontransversal}
\end{align*}
of $\P_4$ satisfy $p\mt q\mt r\mt p$, but $pqr\ne pr$: the sandwich matrix of the rank-$2$ $\D$-class of $\P_4$ has a non-flat triangle, and the $*$-subsemigroup generated by $p,q,r$ is a non-orthodox completely simple regular $*$-semigroup $\mathcal M[\Z_2;3,3;P]$ with $p_{12}\ne1$.
\end{remark}

\begin{remark}[normal $*$-bands]\label{rem:normal}
For a strong semilattice $Y$ of square $*$-bands $A_\alpha\times A_\alpha$ with structure maps $f_{\alpha\beta}\colon A_\alpha\to A_\beta$ ($\alpha\ge\beta$), \eqref{eq:Q1} holds iff all $f_{\alpha\beta}$ are injective: a projection $(c,c)$ of a lower component lies below $(a,b)$ iff $f(a)=f(b)=c$. If $Y$ has a least element $0$, injectivity lets us regard every $A_\alpha$ as a subset of $A=A_0$, and the band becomes the $*$-subsemigroup $\{(\alpha,a,b): a,b\in A_\alpha\}$ of $Y\times(A\times A)$. Since semilattices and square bands embed in partition monoids and $\P_X\times\P_Y\le\P_{X\sqcup Y}$, a finite normal $*$-band embeds in a partition monoid iff it satisfies \eqref{eq:Q1}.
\end{remark}

%==============================================================================
\section{The digraph conditions and the triangle condition}\label{sec:digraph}

\todo{Phase 2. Port from the notes, Section 4: definition of $(\mathrm D_E)$ for admissible digraphs, validity in $\P_X$ (Thm 4.2), the triangle condition as master condition (Thm 4.14, finite case; the infinite case, the unrolled hexagon, as a remark), the witness $\P(Q_0)$ (Thm 4.19), the Mace4 result on order $\le10$.}

%==============================================================================
\section{Wiring conditions}\label{sec:wiring}

\subsection{The shape of the conditions}\label{subsec:shape}

All the conditions of this section are statements about a matched path $a\mt c\mt b$, with the conclusion $a\mt b$, under hypotheses on one, two or three auxiliary projections. In $\P_X$ the hypotheses $a\mt c$, $c\mt b$ pair the $\T$-blocks $A_i\leftrightarrow C_i\leftrightarrow B_i$ through the classes $\Gamma_i$ and $\Delta_i$, but a $\T(a)$-block and a $\T(b)$-block in corresponding classes need not meet or even be connected through $a$ and $b$: the matching does not transport. The conditions say that it does transport as soon as the auxiliaries are wired to the path in one of the following ways, which we name for later reference: an auxiliary $x$ \emph{commutes} with a vertex ($xp=px$); it \emph{lies below} or \emph{above} a vertex; it \emph{sees two vertices alike} ($\theta_x(p)=\theta_x(q)$); it \emph{is seen alike} by two vertices ($\theta_p(x)=\theta_q(x)$); or it is \emph{half-matched} with a vertex ($\theta_p(x)=p$). The conditions are the following; \Wire\ and \Rcond\ are the first two we found, and are consequences of \Qcond\ and \Rp\ respectively (Corollaries~\ref{cor:R} and~\ref{cor:QS}), which is the order in which they are proved below.
\begin{align*}
\Wire:\quad& a\mt c\mt b\mt d,\ cd=dc,\ \theta_a(cd)=\theta_b(cd)\ \Longrightarrow\ a\mt d;\\
\Rcond:\quad& a\mt c\mt b,\ xa=ax,\ xc=cx,\ \theta_x(b)=\theta_x(c)\ \Longrightarrow\ a\mt b;\\
\Vcond:\quad& bcb=b,\ ata\le b,\ ctc\le a\ \Longrightarrow\ ata=btb;\\
\Rp:\quad& a\mt c\mt b,\ xc=cx,\ \theta_x(a)=\theta_x(c),\ ya=ay,\ \theta_x(b)=\theta_y(b)\ \Longrightarrow\ a\mt b;\\
\Pone:\quad& a\mt c\mt b\mt d,\ xa=ax,\ xc=cx,\ \theta_a(x)=\theta_b(x),\ \theta_x(c)=\theta_x(d)\ \Longrightarrow\ a\mt d;\\
\Qcond:\quad& a\mt c\mt b\mt d,\ xa=ax,\ \theta_c(x)=\theta_d(x),\ ya=ay,\ yb=by,\ \theta_y(c)=\theta_y(b)\ \Longrightarrow\ a\mt d;\\
\Kfour:\quad& a\mt c\mt b,\ aya=a,\ b\le z,\ yc=cy,\ \theta_z(y)=\theta_b(a),\ \theta_z(a)=\theta_a(b)\ \Longrightarrow\ a\mt b;\\
\Kzero:\quad& a\mt c\mt b,\ yc=cy,\ \theta_y(b)=\theta_b(a),\ \theta_b(z)=b,\ \theta_z(a)=\theta_z(b),\ \theta_y(z)=\theta_a(b)\ \Longrightarrow\ a\mt b.
\end{align*}
\Wire\ and \Vcond\ look different: in \Wire\ the path has length three and the auxiliary is the product $cd$ of two commuting vertices; \Vcond\ has an order conclusion and a single half-matching hypothesis. But both are of the same kind. \Wire\ is a consequence of \Qcond\ (Corollary~\ref{cor:QS}), and \Vcond\ says that, when $b$ is half-matched with $c$ and the restrictions $\theta_a(t)\le b$ and $\theta_c(t)\le a$ are wired crosswise, the restrictions $\theta_a(t)$ and $\theta_b(t)$ of $t$ coincide; its proof (Theorem~\ref{thm:V}) reduces it by equational reasoning to the statement $\theta_c(\theta_b(t))=\theta_c(t)$, that $c$ sees $t$ and $\theta_b(t)$ alike, which is of the same kind as the hypotheses above. The two conditions \Eone\ and \Esix\ of the census (Section~\ref{subsec:census}) are consequences of \Kzero\ and \Kfour, respectively, together with \eqref{eq:Q1} and the chain condition (W$_2$) (Proposition~\ref{prop:K0E1}).

Table~\ref{tab:witnesses} lists, for each condition, the projection algebra that violates it and satisfies all the others; these are the nine minimal non-embeddable algebras of Theorem~\ref{thm:main-census}. \todo{Table: condition, variables, witness, order, where defined.}

\begin{table}[htbp]
\centering\small
\begin{tabular}{@{}llll@{}}
\toprule
condition & variables & witness & order\\
\midrule
\Wire & $4$ & $\mathfrak P$ & $6$\\
\Rcond & $3+1$ & $\mathfrak P_7$ & $7$\\
\Vcond & $4$ & $\mathfrak P_8^7$ & $8$\\
\Rp & $3+2$ & $\mathfrak P_8^3$ & $8$\\
\Pone & $4+1$ & $\mathfrak P_8^2$ & $8$\\
\Qcond & $4+2$ & $\mathfrak P_8^5$ (and $\mathfrak P_8^4$ via \Ptwo) & $8$\\
\Kfour & $3+2$ & $\mathfrak P_8^6$ (via \Esix) & $8$\\
\Kzero & $3+2$ & $\mathfrak P_8^1$ (via \Eone) & $8$\\
\bottomrule
\end{tabular}
\caption{The wiring conditions and their witnesses.}\label{tab:witnesses}
\end{table}

\subsection{The closing lemma}\label{subsec:closing}

Several proofs end in the same way. For a path $a\mt c\mt b$ in $\P(\P_X)$ consider
\[
\text{(G1)}\quad \varepsilon_c\le\varepsilon_a\vee\varepsilon_b,\qquad\qquad \text{(G2)}\quad\text{no class of }\varepsilon_a\vee\varepsilon_b\vee\varepsilon_c\text{ contains two }\T(c)\text{-blocks.}
\]

\begin{lemma}[closing lemma]\label{lem:closing}
If $a\mt c$, $c\mt b$, \textup{(G1)} and \textup{(G2)} hold in $\P_X$, then $a\mt b$: every class of $\varepsilon_a\vee\varepsilon_b$ meeting $\supp a\cup\supp b\cup\supp c$ contains exactly the $\T$-blocks $A_i,B_i,C_i$ for one $i$.
\end{lemma}

\begin{proof}
By (G1), $\varepsilon_a\vee\varepsilon_b=\varepsilon_a\vee\varepsilon_b\vee\varepsilon_c\ge\varepsilon_a\vee\varepsilon_c,\ \varepsilon_b\vee\varepsilon_c$. Let $D$ be the class of $\varepsilon_a\vee\varepsilon_b$ containing $A_i$. Then $\Gamma_i\subseteq D$, so $C_i\subseteq D$, and $\Delta_i\subseteq D$, so $B_i\subseteq D$. If $A_j\subseteq D$ or $B_j\subseteq D$ for some $j\ne i$, then $C_j\subseteq D$ as well, contradicting (G2). The same applies to the class of any $B_i$ or $C_i$, so the classes meeting $\supp a\cup\supp b$ pair the $A$'s with the $B$'s bijectively, and $a\mt b$ by Lemma~\ref{lem:PX}(c).
\end{proof}

\subsection{The condition \Rp\ and its corollary \Rcond}\label{subsec:Rp}

\begin{theorem}\label{thm:Rp}
\Rp\ holds in every partition monoid: if $a,b,c,x,y$ are projections of $\P_X$ with
\[
a\mt c,\quad c\mt b,\quad xc=cx,\quad \theta_x(a)=\theta_x(c),\quad ya=ay,\quad \theta_x(b)=\theta_y(b),
\]
then $a\mt b$.
\end{theorem}

The proof has two parts. An algebraic part, valid in every regular $*$-semigroup whose projection algebra satisfies \eqref{eq:Q1}, derives a number of identities from the hypotheses and reduces the theorem to the case $x\le y$; the combinatorial part is a chain argument in $\P_X$ which shows that no class of $\varepsilon_a\vee\varepsilon_b\vee\varepsilon_c$ contains two $\T(c)$-blocks, followed by a count. Most of the algebraic part was found with Prover9, and the combinatorial statements were found by testing candidates against all instances of the hypotheses in $\P(\P_4)$.

\begin{proof}
Throughout, $r=\theta_x(b)=\theta_y(b)$, $f=xc=cx$ and $e=ya=ay$.

\smallskip
\emph{Step 1: algebraic consequences.} The identities of this step hold in every regular $*$-semigroup satisfying \eqref{eq:Q1}, under the hypotheses of the theorem.

(i) $r\le x$ and $r\le y$, so $xy=yx=:m$ by \eqref{eq:Q1}, and $r\le m$.

(ii) $\theta_x(c)=xcx=xc=f$, so the hypothesis $\theta_x(a)=\theta_x(c)$ reads $xax=f$. Hence $xa=xa\,ax\,xa=(xax)a=fa$ and, applying $*$, $ax=af$.

(iii) \emph{$f\le y$ and $cy=yc=f$.} First, $\theta_a(r)=a(yby)a=(ay)b(ya)=ebe$. Next, for every element $s$, $\theta_x(\theta_e(s))=x\,ya\,s\,ay\,x=m\,asa\,m=\theta_m(\theta_a(s))$, since $xya=ma$ and $ayx=am$. Applied to $s=b$, with $\theta_a(r)=ebe$: $\theta_m(aba)=\theta_x(ebe)=\theta_x(\theta_a(r))=xa\,xbx\,ax=(xax)b(xax)=fbf$, and $fbf=f$ because $f\le c$ and $cbc=c$ (for $q\le c$, $qbq=qc\,b\,cq=q(cbc)q=q$). So $f=\theta_m(aba)\le m\le y$. Now $f$ is a common lower bound of $c$ and $y$, so $cy=yc$ by \eqref{eq:Q1}; the projection $q=cy$ lies below $c$, so $qbq=q$ by the identity just used, and $qbq=c(yby)c=crc=c(xbx)c=(cx)b(xc)=fbf=f$. Thus $cy=yc=f$.

(iv) \emph{$\theta_a(x)=\theta_a(y)=e$ and $e\mt f$.} By (ii) and (iii), $axa=afa=a\,cy\,a=ac\,ya=aca\,y=ay=e$. Using $ay=ya$, $cy=yc$, $aca=a$ and $cac=c$: $efe=ya\,yc\,ya=y\,ac\,ya=y\,aca\,y=e$ and $fef=yc\,ya\,yc=y\,cac\,y=f$.

(v) \emph{$yax=ax$, $xya=xa$, $axb=yab$ and $xab=xcb$.} We have $cbf=cb\,cx=(cbc)x=cx=f$ and, by (ii), $ax=af=acx$; hence $axbf=afbf=a(fbf)=af=ax$. Since $r\le y$, $y\,xbx=yr=r=xbx$; therefore $yax=y\,axbf=y\,a(xbx)c=a\,y(xbx)\,c=a(xbx)c=axbf=ax$. Applying $*$ gives $xay=xa$, i.e.\ $xya=xa$. Next $xb=(xbx)b=rb=(yby)b=yb$, so $axb=ayb=yab$. Finally $xcb=(xax)b=x(axb)=x(yab)=(xya)b=xab$.

(vi) \emph{$y\,ab=y\,acb$.} Multiplying $xab=xcb$ on the left by $a$ and using $axa=e$ and $axc=afc=af$ (as $f\le c$): $e\,b=af\,b$; and $y\,acb=(ya)(cb)=a(yc)b=afb$. So $eb=y\,ab=y\,acb=ecb$.

(vii) \emph{Reduction to $x\le y$.} The projection $m=xy$ satisfies the hypotheses in place of $x$: $mc=xyc=xf=f=fy=cxy=cm$, as $f\le x$ and $f\le y$; $\theta_m(c)=mcm=f$, and $\theta_m(a)=xy\,a\,yx=xex=x(afa)x=(xa)f(ax)=(fa)f(af)=f(afa)f=fef=f$ by (ii) and (iv); and $\theta_m(b)=\theta_m(\theta_x(b))=\theta_m(r)=r=\theta_y(b)$, since $m\le x$ and $r\le m$. As the conclusion does not involve $x$, we assume from now on that $x\le y$.

\smallskip
\emph{Step 2: the setting in $\P_X$.} We use the notation of Section~\ref{subsec:PX}: $A_i$, $C_i$, $B_i$, $\Gamma_i$, $\Delta_i$; for a block $E$ of $c$ we write $\Gamma(E)$ and $\Delta(E)$ for its classes in $\varepsilon_a\vee\varepsilon_c$ and $\varepsilon_b\vee\varepsilon_c$. Put $Y=X\setminus\supp y$. We shall show that
\begin{equation}\label{eq:Rp-goal}
\text{no class of }\varepsilon_a\vee\varepsilon_b\vee\varepsilon_c\text{ contains two }\T(c)\text{-blocks,}
\end{equation}
and then conclude by a count. Since $A_i,B_i,C_i$ lie in one class of $\varepsilon_a\vee\varepsilon_b\vee\varepsilon_c$, \eqref{eq:Rp-goal} implies that every class of $\varepsilon_a\vee\varepsilon_b$ contains at most one $\T(a)$-block and at most one $\T(b)$-block, and that a class containing $A_i$ contains no $\T(b)$-block other than $B_i$.

By Lemma~\ref{lem:PX}(b) applied to the commuting pairs $y,a$ and $y,c$ and $x,c$: every block of $a$, of $c$ and of $x$ lies inside $\supp y$ or is an $\N(y)$-block; an $\N(y)$-block is an $\N(a)$-block or a union of $\T(a)$-blocks, and an $\N(c)$-block or a union of $\T(c)$-blocks; every $\N(c)$-block lies inside $\supp x$ (and is then a union of $\T(x)$-blocks) or is an $\N(x)$-block. Since $x\le y$, every block of $x$ is an $\N(y)$-block or a union of $\T(y)$-blocks (Lemma~\ref{lem:order}(a)). The classes of $\varepsilon_y\vee\varepsilon_a$ meeting $\supp y$ are the $\T(e)$-blocks $U_1,\dots,U_t$ (unions of $\T(a)$-blocks and of $\T(y)$-blocks) and the $\N(a)$-blocks inside $\supp y$; the classes of $\varepsilon_y\vee\varepsilon_c$ meeting $\supp y$ are the $\T(f)$-blocks $V_1,\dots,V_{t'}$ and the $\N(c)$-blocks inside $\supp y$.

\emph{Step 3: a second reduction.} Let $M$ be an $\N(y)$-block which is a union of $\T(a)$-blocks. If $M$ were an $\N(c)$-block, it would be a class of $\varepsilon_a\vee\varepsilon_c$ by itself, containing $\T(a)$-blocks but no $\T(c)$-block, contradicting $a\mt c$; so $M$ is a union of $\T(c)$-blocks, the class of any $A_i\subseteq M$ lies inside $M$ and consists of $\T$-blocks only, and $\Gamma_i=A_i=C_i$. Symmetrically, an $\N(y)$-block which is a union of $\T(c)$-blocks is a union of $\T(a)$-blocks. Let $x',y'$ be obtained from $x,y$ by declaring all these blocks $M$ transversal. We claim that $a,b,c,x',y'$ still satisfy the hypotheses, with $x'\le y'$; as the conclusion does not change, we may then assume that \emph{every $\N(y)$-block is an $\N(a)$-block and an $\N(c)$-block}, so that all $A_i$ and all $C_i$ lie inside $\supp y$. Indeed, the partitions $\varepsilon_{x'}=\varepsilon_x$, $\varepsilon_{y'}=\varepsilon_y$ are unchanged, and each $M$ is a class by itself in $\varepsilon_x\vee\varepsilon_a$, $\varepsilon_x\vee\varepsilon_c$, $\varepsilon_y\vee\varepsilon_a$ and $\varepsilon_y\vee\varepsilon_c$, being a block of $x$ and of $y$ which is a union of blocks of $a$ and of $c$. Hence $x'c=cx'$, $y'a=ay'$ and $y'c=cy'$ by the converse of Lemma~\ref{lem:PX}(b) (if every class of $\varepsilon_p\vee\varepsilon_q$ is of one of the four types, the description of $pq$ in the proof of Lemma~\ref{lem:PX} shows that $pq$ is symmetric, so $pq=qp$): the class $M$ is now of type (i), the others are unchanged. Further $\theta_{x'}(a)=\theta_{x'}(c)$, because the blocks of these projections are those of $\theta_x(a)=\theta_x(c)$ except that each $M$ becomes a transversal block of both; and $\theta_{x'}(b)=\theta_{y'}(b)$ will follow from Step~4(c) below, whose proof uses only $x\le y$ and $\theta_x(b)=\theta_y(b)$.

\emph{Step 4: the classes of $\varepsilon_y\vee\varepsilon_b$.} Let $\sigma=\varepsilon_y\vee\varepsilon_b$.

(a) \emph{$\varepsilon_x\vee\varepsilon_b=\sigma$.} Let $W$ be a block of $x$ which is a union of $\T(y)$-blocks. If $W$ is a $\T(x)$-block, let $\tilde K$ be its class in $\varepsilon_x\vee\varepsilon_b$; then $\tilde K\cap\supp x$ is a block of $\theta_x(b)=\theta_y(b)$ contained in $\supp y$, hence equal to $K\cap\supp y$ for a class $K$ of $\sigma$, and $W\subseteq K$. If $W$ is an $\N(x)$-block, it is itself a block of $\theta_x(b)=\theta_y(b)$ contained in $\supp y$, so again $W=K\cap\supp y$ for a class $K$ of $\sigma$. The remaining blocks of $x$ are $\N(y)$-blocks. So every block of $x$ lies in a single class of $\sigma$, i.e.\ $\varepsilon_x\le\sigma$, and $\varepsilon_x\vee\varepsilon_b=\sigma$.

(b) \emph{Two kinds of classes.} By (a) and the two descriptions of $r$, for every class $K$ of $\sigma$ meeting $\supp y$ either $K\cap\supp y=\tilde N$ is an $\N(x)$-block (a union of $\T(y)$-blocks), and then $K\cap\supp b=\emptyset$ since $\tilde N$ is a non-transversal block of $\theta_x(b)$; or $K\cap\supp y=K\cap\supp x$. We call the $\N(x)$-blocks inside $\supp y$ the \emph{$\tilde N$-blocks}; each $\tilde N$-block is $K\cap\supp y$ for a class $K$ of $\sigma$ with $K\cap\supp b=\emptyset$, and every other class meeting $\supp y$ has $K\cap\supp y\subseteq\supp x$.

(c) In the situation of Step~3, the blocks of $\theta_{x'}(b)$ and $\theta_{y'}(b)$ are indexed by the same classes $K$ of $\sigma$ (by (a)), and they can differ only for a class $K\supseteq M$, where they are $(K\cap\supp x)\cup M$ and $(K\cap\supp y)\cup M$, both transversal iff $K$ meets $\supp b$. If these differ, then by (b) $K\cap\supp y$ is an $\tilde N$-block and $K\cap\supp b=\emptyset$; but then, for $A_i=C_i\subseteq M$, every block of $b$ meeting $K$ lies in $K$, and every block of $c$ meeting a block $Z$ of $b$ inside $K$ lies in $K$ (a block of $c$ in $Y$ is an $\N(y)$-block, hence lies in a single class of $\sigma$; a block of $c$ inside $\supp y$ meets $Z\cap\supp y\subseteq\tilde N$, so it is $\tilde N$ or contained in it), so a chain from $C_i\subseteq K$ to $B_i$ in $\Delta_i$ stays inside $K$, contradicting $K\cap\supp b=\emptyset$. This completes Step~3.

(d) \emph{Every $\tilde N$-block is an $\N(c)$-block, and its class $\Delta(\tilde N)$ contains no $\T(b)$-block and satisfies $\Delta(\tilde N)\cap\supp y=\tilde N$.} Let $K$ be the class of $\sigma$ with $K\cap\supp y=\tilde N$. As in (c), every block of $b$ or of $c$ reachable by a chain in $\varepsilon_b\vee\varepsilon_c$ from a block of $b$ inside $K$ lies inside $K$. If $\tilde N$ were a union of $\T(c)$-blocks (the alternative, as $x$ and $c$ commute), a block of $b$ meeting some $C_j\subseteq\tilde N$ would lie in $K$, and then $B_j\subseteq\Delta_j\subseteq K$, contradicting $K\cap\supp b=\emptyset$. So $\tilde N$ is an $\N(c)$-block, $\Delta(\tilde N)\subseteq K$, and the claims follow.

\emph{Step 5: the classes of $\varepsilon_x\vee\varepsilon_a$.} By Lemma~\ref{lem:PX}(a), the blocks of $\theta_x(a)$ inside $\supp x$ are the sets $K\cap\supp x$ for the classes $K$ of $\varepsilon_x\vee\varepsilon_a$ meeting $\supp x$, transversal iff $K$ meets $\supp a$; those of $\theta_x(c)$ are the $\T(f)$-blocks $V_k$ (transversal) and the $\N(c)$-blocks $\tilde N'$ inside $\supp x$ (non-transversal; they are the classes of $\varepsilon_x\vee\varepsilon_c$ of type (ii)). Since $\theta_x(a)=\theta_x(c)$: for every $k$ there is a class $K_k$ of $\varepsilon_x\vee\varepsilon_a$ with $K_k\cap\supp x=V_k$, and for every $\N(c)$-block $\tilde N'\subseteq\supp x$ a class $K_{\tilde N'}$ with $K_{\tilde N'}\cap\supp x=\tilde N'$ and $K_{\tilde N'}\cap\supp a=\emptyset$. Consequently \emph{every block of $c$ lies inside a single class of $\varepsilon_x\vee\varepsilon_a$}: a $\T(c)$-block lies in some $V_k\subseteq K_k$; an $\N(c)$-block inside $\supp x$ is some $\tilde N'\subseteq K_{\tilde N'}$; and the other $\N(c)$-blocks are $\N(x)$-blocks. Hence $\varepsilon_a\vee\varepsilon_c\le\varepsilon_x\vee\varepsilon_a$, and for an $\N(c)$-block $\tilde N'\subseteq\supp x$ we get $\Gamma(\tilde N')\subseteq K_{\tilde N'}$; so \emph{$\Gamma(\tilde N')$ contains no $\T(a)$-block, hence no $\T(c)$-block, and $\tilde N'$ is the only block of $c$ inside $\supp x$ that it contains} (a block of $c$ in $\Gamma(\tilde N')$ inside $\supp x$ lies in $K_{\tilde N'}\cap\supp x=\tilde N'$).

\emph{Step 6: the four kinds of blocks of $c$ and their classes.} By Steps~2--4, every block of $c$ is of exactly one of the following kinds: a $\T(c)$-block $C_i$ (inside $\supp y$, by Step~3); an $\N(c)$-block inside $\supp x$; an $\tilde N$-block; or an $\N(y)$-block inside $Y$, which is an $\N(a)$-block by Step~3 and is therefore a class of $\varepsilon_a\vee\varepsilon_c$ by itself. Collecting Steps~4(d) and 5:
\begin{itemize}[nosep]
\item[(1)] the blocks of $c$ in $\Gamma_i$ are $C_i$ and $\tilde N$-blocks (an $\N(c)$-block inside $\supp x$ lies in a class without $\T(a)$-blocks, and an $\N(y)$-block in $Y$ is a class by itself);
\item[(2)] the blocks of $c$ in $\Delta_i$ are $C_i$, $\N(c)$-blocks inside $\supp x$ and $\N(y)$-blocks in $Y$ (an $\tilde N$-block lies in a class without $\T(b)$-blocks);
\item[(3)] for an $\N(c)$-block $\tilde N'\subseteq\supp x$, the blocks of $c$ in $\Gamma(\tilde N')$ are $\tilde N'$ and $\tilde N$-blocks, and $\Gamma(\tilde N')$ contains no $\T(a)$-block;
\item[(4)] for an $\tilde N$-block $\tilde N$, the blocks of $c$ in $\Delta(\tilde N)$ are $\tilde N$ and $\N(y)$-blocks in $Y$ (as $\Delta(\tilde N)\cap\supp y=\tilde N$), and $\Delta(\tilde N)$ contains no $\T(b)$-block;
\item[(5)] an $\N(y)$-block $E\subseteq Y$ is a class of $\varepsilon_a\vee\varepsilon_c$ by itself.
\end{itemize}

\emph{Step 7: proof of \eqref{eq:Rp-goal}.} Two blocks $E,E'$ of $c$ lie in one class of $\varepsilon_a\vee\varepsilon_b\vee\varepsilon_c$ iff they are joined by a sequence $E=E_0,E_1,\dots,E_n=E'$ of blocks of $c$ in which consecutive blocks lie in a common class of $\varepsilon_a\vee\varepsilon_c$ or in a common class of $\varepsilon_b\vee\varepsilon_c$: if two blocks $W,W'$ (of $a$, $b$ or $c$) meet at a point $p$, and $E(p)$ is the block of $c$ containing $p$, then $W$ and $E(p)$ lie in a common class of $\varepsilon_a\vee\varepsilon_c$ or of $\varepsilon_b\vee\varepsilon_c$ (according to whether $W$ is a block of $a$ or of $b$; if $W$ is a block of $c$ then $W=E(p)$), so a chain of blocks meeting consecutively can be replaced by the sequence of the blocks of $c$ at its meeting points. Now start such a sequence at $E_0=C_p$ and use (1)--(5) to see which blocks of $c$ can occur. From $C_p$ one can move, within $\Gamma_p$, only to $\tilde N$-blocks of $\Gamma_p$, and within $\Delta_p$ only to $\N(c)$-blocks $\tilde N'\subseteq\supp x$ of $\Delta_p$ and to $\N(y)$-blocks $E\subseteq Y$ of $\Delta_p$. From an $\tilde N$-block $\tilde N$ of $\Gamma_p$ one can move within $\Gamma(\tilde N)=\Gamma_p$ (back) or within $\Delta(\tilde N)$ to $\N(y)$-blocks $E\subseteq Y$; from such an $E$ one can move within $\Delta(E)=\Delta(\tilde N)$ only, and $\Gamma(E)=\{E\}$. From $\tilde N'\subseteq\supp x$ in $\Delta_p$ one can move within $\Delta(\tilde N')=\Delta_p$ (back) or within $\Gamma(\tilde N')$ to $\tilde N$-blocks of $\Gamma(\tilde N')$, from which one can move within $\Gamma(\tilde N')$ (back) or within their $\Delta$-classes to $\N(y)$-blocks in $Y$, which are dead ends as before. From an $\N(y)$-block $E\subseteq Y$ in $\Delta_p$ one can move within $\Delta_p$ only. So the blocks of $c$ reachable from $C_p$ are: $C_p$; $\tilde N$-blocks (of $\Gamma_p$, or of the classes $\Gamma(\tilde N')$ for $\tilde N'\subseteq\supp x$ in $\Delta_p$); $\N(c)$-blocks inside $\supp x$ lying in $\Delta_p$; and $\N(y)$-blocks in $Y$. None of these is a $\T(c)$-block other than $C_p$, which proves \eqref{eq:Rp-goal}.

\emph{Step 8: counting.} Let $l\in\{1,\dots,t\}$ and $I_l=\{p:A_p\subseteq U_l\}$; by Step~3 every $A_p$ lies in some $U_l$. Consider the class of $U_l$ in $\varepsilon_e\vee\varepsilon_f$ (where $\varepsilon_e=\varepsilon_y\vee\varepsilon_a$, $\varepsilon_f=\varepsilon_y\vee\varepsilon_c$); by $e\mt f$ (Step~1(iv)) it contains exactly one $\T(f)$-block, $V_{k(l)}$ say. Every $\Gamma_p$ is contained in a class of $\varepsilon_e\vee\varepsilon_f$ (its blocks of $a$ and of $c$ are contained in blocks of $e$ and of $f$), so $C_p\subseteq V_{k(l)}$ for $p\in I_l$, and conversely $C_q\subseteq V_{k(l)}$ implies $A_q\subseteq U_l$; thus $I_l=\{q:C_q\subseteq V_{k(l)}\}$. Now compare the two sides of $y\,ab=y\,acb$ (Step~1(vi)). The upper blocks of $ab$ are the $\N(a)$-blocks and the sets $D\cap\supp a$ for the classes $D$ of $\varepsilon_a\vee\varepsilon_b$ meeting $\supp a$, the latter being transversal iff $D$ meets $\supp b$, with lower part $D\cap\supp b$. By \eqref{eq:Rp-goal}, $D\cap\supp a$ is a single $\T(a)$-block; write $D(A_p)$ for the class of $A_p$. In the product $y\cdot(ab)$ the middle join is $\varepsilon_y\vee\varepsilon_{ab}$, where $\varepsilon_{ab}$ is the upper partition of $ab$; the class of $U_l$ in it is $U_l$ itself, because $U_l$ is a union of $\T(y)$-blocks and of the blocks $A_p$ ($p\in I_l$) of $\varepsilon_{ab}$, and no $\N(a)$-block or $\N(y)$-block meets it. So $y\,ab$ has the upper block $U_l$, with lower part $\bigcup_{p\in I_l}(D(A_p)\cap\supp b)$ (transversal iff this is nonempty). In $y\,acb=(ya)(cb)=e\,(cb)$ the middle join is $\varepsilon_e\vee\varepsilon_c$, whose class containing $U_l$ contains no other $\T(e)$-block and, being contained in the class of $U_l$ in $\varepsilon_e\vee\varepsilon_f$, contains exactly the $\T(c)$-blocks $C_p$ with $p\in I_l$ (all of which it does contain, as $\Gamma_p$ is a chain in $\varepsilon_e\vee\varepsilon_c$); so $e\,cb$ has the upper block $U_l$ with lower part $\bigcup_{p\in I_l}B_p$. Hence
\[
\bigcup_{p\in I_l}\bigl(D(A_p)\cap\supp b\bigr)=\bigcup_{p\in I_l}B_p .
\]
By \eqref{eq:Rp-goal}, $D(A_p)\cap\supp b\subseteq B_p$ for each $p$, and the $B_p$ are pairwise disjoint and nonempty; so $D(A_p)\cap\supp b=B_p$, i.e.\ $B_p\subseteq D(A_p)$, for every $p\in I_l$ and every $l$, that is, for every $p$.

\emph{Conclusion.} Every class of $\varepsilon_a\vee\varepsilon_b$ meeting $\supp a$ is some $D(A_p)$, which contains exactly one $\T(a)$-block and exactly one $\T(b)$-block ($B_p$) by \eqref{eq:Rp-goal} and Step~8; and every class meeting $\supp b$ contains some $B_p$, hence is $D(A_p)$. By Lemma~\ref{lem:PX}(c), $a\mt b$.
\end{proof}

\begin{remark}\label{rem:Rp}
The conclusion obtained is stronger than $a\mt b$: by \eqref{eq:Rp-goal} and Step~8 the classes of $\varepsilon_a\vee\varepsilon_b\vee\varepsilon_c$ meeting $\supp a\cup\supp b\cup\supp c$ are exactly the classes containing the triples $A_i,B_i,C_i$, so that the three matchings are coherent and $ab=acb$. The count of Step~8 alone cannot replace \eqref{eq:Rp-goal}: it is blind to the merging of two $\T(a)$-blocks of the same $U_l$ into one class of $\varepsilon_a\vee\varepsilon_b$, which is exactly what happens in the witness $\mathfrak P_8^3$.
\end{remark}

\begin{corollary}\label{cor:R}
\Rcond\ holds in every projection algebra satisfying \Rp; in particular in every $\P(\P_X)$.
\end{corollary}

\begin{proof}
Given the hypotheses of \Rcond, put $x'=xc=cx$ and $y'=x$. Then $x'\le c$, so $x'c=cx'$; $\theta_{x'}(a)=cx\,a\,xc=c(xax)c=c(xa)c=x(cac)=xc=\theta_{x'}(c)$, using $xa=ax$, $xc=cx$ and $cac=c$; $y'=x$ commutes with $a$; and $\theta_{x'}(b)=c\,\theta_x(b)\,c=c\,\theta_x(c)\,c=cxc=xc=\theta_x(c)=\theta_x(b)=\theta_{y'}(b)$. So \Rp\ gives $a\mt b$. (The substitution was found by Prover9.)
\end{proof}

\subsection{The condition \Vcond}\label{subsec:V}

\begin{theorem}\label{thm:V}
\Vcond\ holds in every partition monoid: if $a,b,c,t$ are projections of $\P_X$ with $bcb=b$, $ata\le b$ and $ctc\le a$, then $ata=btb$.
\end{theorem}

\begin{proof}
Write $z=ata$, $w=ctc$, $s=btb$. We first reduce the claim, by computations valid in every projection algebra satisfying \eqref{eq:Q1}, to a single statement about $\P_X$.

\smallskip
\emph{Step 1: three identities.} (i) \emph{If $k\le a$ and $ab=ba$ then $bkb\le a$}: $a(bkb)a=(ab)k(ba)=(ba)k(ab)=b(aka)b=bkb$. (ii) \emph{If $ata\le b$ then $ab=ba$ and $ata=ete$, where $e=ab$}: $z=ata$ lies below $a$ and below $b$, so $a$ and $b$ commute by \eqref{eq:Q1}; $e=ab$ is a projection with $e\le a$, $e\le b$, and $eze=a(bzb)a=aza=z$, so $z\le e$ and $z=eze=e(ata)e=(ea)t(ae)=ete$. (iii) \emph{If $bcb=b$ then $s=b(csc)b$}: $s=btb=(bcb)t(bcb)=b\,c(btb)c\,b$.

\emph{Step 2: reduction.} Suppose we know that
\begin{equation}\label{eq:lemD}
csc=w,\qquad\text{i.e.}\qquad \theta_c(\theta_b(t))=\theta_c(t) .
\end{equation}
Then $s=b(csc)b=bwb$ by (iii), and $bwb\le a$ by (i), since $w\le a$ and $ab=ba$ (by (ii), as $z\le b$). So $s\le a$ and $s\le b$, hence $s\le e=ab$ and $s=ese=e(btb)e=(eb)t(be)=ete$; and $z=ete$ by (ii). Thus $s=z$. It remains to prove \eqref{eq:lemD} in $\P_X$; in fact we show that $\theta_c(s)$ and $\theta_c(t)$ have the same blocks.

\smallskip
\emph{Step 3: the setting in $\P_X$.} Put $J_a=\varepsilon_a\vee\varepsilon_t$, $J_b=\varepsilon_b\vee\varepsilon_t$, $J_c=\varepsilon_c\vee\varepsilon_t$. By Lemma~\ref{lem:PX}(a) the blocks of $z$ are the $\N(a)$-blocks and the sets $Z_K=K\cap\supp a$ for the classes $K$ of $J_a$ meeting $\supp a$, $Z_K$ being transversal iff $K$ meets $\supp t$; similarly the blocks of $w$ are the $\N(c)$-blocks and the sets $W_L=L\cap\supp c$ ($L\in J_c$), and those of $s$ are the $\N(b)$-blocks and the sets $S_H=H\cap\supp b$ ($H\in J_b$). From $w\le a$ and Lemma~\ref{lem:order}(a): every block of $a$ lies in a block of $w$; every block of $w$ is an $\N(a)$-block or a union of $\T(a)$-blocks; and every $\N(a)$-block is an $\N(w)$-block, i.e.\ an $\N(c)$-block or a set $W_L$ with $L\cap\supp t=\emptyset$. From $z\le b$: every block of $b$ lies in a block of $z$; every block of $z$ is an $\N(b)$-block or a union of $\T(b)$-blocks; and every $\N(b)$-block is a block of $z$, non-transversal. From $bcb=b$ and Lemma~\ref{lem:order}(b): every block of $c$ meets at most one $\T(b)$-block, and for every $\T(b)$-block $B$ the class $\Delta(B)$ of $\varepsilon_b\vee\varepsilon_c$ containing $B$ meets $\supp c$. Since blocks of $a$ lie in blocks of $w$, which lie in classes of $J_c$, we have $J_a\le J_c$ (every class of $J_a$ lies in a class of $J_c$); likewise $J_b\le J_a$. For a class $K$ of $J_a$ let $\lambda(K)$ be the class of $J_c$ containing it. Call a block of $c$ \emph{inner} if it meets $\supp a$ and \emph{outer} otherwise. An inner $\T(c)$-block lies in a set $W_L$ which is a union of $\T(a)$-blocks; an inner $\N(c)$-block is a union of $\T(a)$-blocks (being a block of $w$ that is not an $\N(a)$-block); an outer block of $c$ inside a class $L$ of $J_c$ with $W_L\ne\emptyset$ is either a $\T(c)$-block inside $W_L$, $W_L$ being then an $\N(a)$-block, or an $\N(c)$-block which is an $\N(a)$-block.

\emph{Step 4: two facts.} (1) \emph{No $\T(b)$-block lies in a block $M$ which is both an $\N(a)$-block and an $\N(c)$-block.} Otherwise let $B\subseteq M$ be one. The only block of $c$ meeting $B$ is $M$, and the blocks of $b$ meeting $M$ lie inside $M$ ($M$ is a block of $z$, and blocks of $b$ lie in blocks of $z$); so $\Delta(B)=M$ does not meet $\supp c$, contradicting Lemma~\ref{lem:order}(b). Consequently every block which is both an $\N(a)$-block and an $\N(c)$-block is an $\N(b)$-block (a block of $z$ containing no $\T(b)$-block).

(2) \emph{Let $K$ be a class of $J_a$ and $M\subseteq K$ an $\N(a)$-block containing a $\T(b)$-block. Then $M$ is the only block of $z$ inside $K$ that meets $\supp b$; in particular $Z_K$, if nonempty, is an $\N(b)$-block.} By (1) $M$ is not an $\N(c)$-block, so $M=W_L$ with $L=\lambda(K)$, $L\cap\supp t=\emptyset$; this already shows that no other $\N(a)$-block in $K$ contains a $\T(b)$-block (it would also equal $W_{\lambda(K)}$). Suppose $Z_K$ contains a $\T(b)$-block $B$. Every $\T(a)$-block in $K$ lies in $L$ but not in $W_L=M$ (which is disjoint from $\supp a$), hence in an $\N(c)$-block, which is then inner, a union of $\T(a)$-blocks, and contained in $L$. Take a chain $B=B_0,C_1,B_1,C_2,\dots$ from $B$ to a $\T(c)$-block in $\Delta(B)$ (it exists by Lemma~\ref{lem:order}(b)). $C_1$ meets $B\subseteq Z_K\subseteq\supp a$ in a point of a $\T(a)$-block of $K$, so $C_1$ is an inner $\N(c)$-block, $C_1\subseteq L\cap\supp a$; then $B_1$ meets $\supp a$ in a point of a $\T(a)$-block $A_1\subseteq C_1\subseteq L$, whose class $K_1$ satisfies $\lambda(K_1)=L$, and $B_1\subseteq Z_{K_1}$; so $C_2$ meets a $\T(a)$-block of $K_1$, which lies in $L$ but not in $W_L$, and $C_2$ is again an inner $\N(c)$-block. Inductively no block of the chain is a $\T(c)$-block, a contradiction.

\emph{Step 5: the blocks of $\theta_c(s)$ and $\theta_c(t)$ inside $\supp c$.} By Lemma~\ref{lem:PX}(a) the blocks of $\theta_c(s)$ are the $\N(c)$-blocks and the sets $\Theta\cap\supp c$ for the classes $\Theta$ of $\varepsilon_c\vee\varepsilon_s$ meeting $\supp c$, transversal iff $\Theta$ meets $\supp s$; the blocks of $\theta_c(t)=w$ are the $\N(c)$-blocks and the $W_L$, transversal iff $L$ meets $\supp t$. Every block of $s$ lies in a class of $J_c$: an $\N(b)$-block is a block of $z$, hence lies in a class of $J_a$, and $S_H\subseteq H$, a class of $J_b\le J_a\le J_c$. Hence $\varepsilon_c\vee\varepsilon_s\le J_c$: every class $\Theta$ lies in a class $L$ of $J_c$, $\Theta\cap\supp c\subseteq W_L$, and if $\Theta$ meets $\supp s$ then $L$ meets $\supp t$ (a block $S_H\subseteq\Theta$ is transversal only if $H\subseteq L$ meets $\supp t$). So \eqref{eq:lemD} follows once we show, for every class $L$ of $J_c$ with $W_L\ne\emptyset$: \emph{(A) any two $\T(c)$-blocks in $W_L$ lie in a common class of $\varepsilon_c\vee\varepsilon_s$, and (B) if $L$ meets $\supp t$ then this class meets $\supp s$.} Fix $L$; then $W_L$ is a union of $\T(a)$-blocks (\emph{type T}) or an $\N(a)$-block $M$ (\emph{type N}, and then $L\cap\supp t=\emptyset$, so (B) is vacuous).

\emph{Type T.} Let $C,C'$ be $\T(c)$-blocks in $W_L$ and take a chain $C=C_0,T_1,C_1,\dots,T_r,C_r=C'$ in $J_c$ (blocks $C_i$ of $c$, blocks $T_i$ of $t$, all inside $L$). We claim that consecutive inner blocks of $c$ on the chain, say $C_i$ and $C_j$ with $C_{i+1},\dots,C_{j-1}$ outer, lie in a common class of $\varepsilon_c\vee\varepsilon_s$; since $C_0$ and $C_r$ are inner ($W_L\subseteq\supp a$), (A) follows. The outer blocks $C_{i+1},\dots,C_{j-1}$ are $\N(c)$-blocks which are $\N(a)$-blocks (the $\T(c)$-blocks of $L$ lie in $W_L\subseteq\supp a$), hence $\N(b)$-blocks by (1); the blocks $T_{i+1},C_{i+1},\dots,T_j$ lie in a single class $K''$ of $J_a$ (each lies in a class, and consecutive ones meet). Let $x\in T_{i+1}\cap C_i$ and $y\in T_j\cap C_j$; both lie in $\supp a$ ($C_i,C_j$ are inner: a $\T(c)$-block inside $W_L$, or an $\N(c)$-block which is a union of $\T(a)$-blocks), so $x,y\in Z_{K''}$. If $Z_{K''}$ is an $\N(b)$-block, it is a block of $s$ meeting $C_i$ and $C_j$. Otherwise $Z_{K''}$ is a union of $\T(b)$-blocks; let $B_x\ni x$, $B_y\ni y$ be the $\T(b)$-blocks; then $B_x,T_{i+1},C_{i+1},T_{i+2},\dots,T_j,B_y$ is a chain in $J_b$, so $B_x,B_y\subseteq H$ for one class $H$ of $J_b$, and the block $S_H\supseteq B_x\cup B_y$ of $s$ meets $C_i$ and $C_j$. This proves the claim. For (B), let $L$ meet $\supp t$, and take a chain $C=C_0,T_1,C_1,\dots,C_q,T_{q+1}$ in $L$ from $C$ to a $\T(t)$-block $T_{q+1}$. Let $C_j$ ($j\le q$) be the last inner block of $c$ on it; then $C_{j+1},\dots,C_q$ are $\N(b)$-blocks as above, and $T_{j+1},C_{j+1},\dots,T_{q+1}$ lie in one class $K'$ of $J_a$, which meets $\supp t$. The point $x\in T_{j+1}\cap C_j$ lies in $Z_{K'}$, a transversal block of $z$, hence a union of $\T(b)$-blocks (an $\N(b)$-block is a non-transversal block of $z$); the $\T(b)$-block $B_x\ni x$ is joined to $T_{q+1}$ by the chain $B_x,T_{j+1},C_{j+1},\dots,C_q,T_{q+1}$ in $J_b$, so its class $H$ meets $\supp t$, $S_H$ is a transversal block of $s$, and $S_H$ meets $C_j$, which lies in the class of $C$ by the claim.

\emph{Type N.} Here $W_L=M$ is an $\N(a)$-block and a block of $z$. If $M$ is an $\N(b)$-block, it is a block of $s$ containing every $\T(c)$-block of $W_L$, and (A) holds. Otherwise $M$ is a union of $\T(b)$-blocks $B_1,\dots,B_k$; each $\T(c)$-block $C\subseteq M$ lies inside one $B_i$ (it is covered by them and meets at most one), so (A) follows if we show that $B_1,\dots,B_k$ lie in a single class $H$ of $J_b$ (then $S_H\supseteq M$). Let $K_0$ be the class of $J_a$ containing $M$. By (2), $M$ is the only block of $z$ inside $K_0$ meeting $\supp b$; so every $\N(a)$-block in $K_0$ other than $M$ is an $\N(b)$-block, and $Z_{K_0}$ is an $\N(b)$-block if nonempty. Hence every block of $b$ inside $K_0$ is one of: a $\T(b)$-block $B_i$; an $\N(a)$-block $M'\ne M$ in $K_0$; or $Z_{K_0}$; and for a point $p\in K_0$ we write $\beta(p)$ for the block of $b$ containing it, which is $Z_{K_0}$ whenever $p\in\supp a$. Let $C\subseteq B_1$ and $C'\subseteq B_2$ be $\T(c)$-blocks and take a chain $C=C_0,T_1,C_1,\dots,T_r,C_r=C'$ in $L$. Every block of $t$ meeting $M$, $Z_{K_0}$ or an $\N(a)$-block in $K_0$ lies in $K_0$; in particular $T_1,T_r\subseteq K_0$. Consider a block $C_l$ of the chain with $T_l\subseteq K_0$, and let $u\in T_l\cap C_l$. If $C_l$ is a $\T(c)$-block, then $C_l\subseteq M$ lies inside one $B_i$, and $\beta(u)=\beta(v)=B_i$ for the point $v\in T_{l+1}\cap C_l$. If $C_l$ is an outer $\N(c)$-block, it is an $\N(a)$-block inside $K_0$ and an $\N(b)$-block by (1), and $\beta(u)=\beta(v)=C_l$. If $C_l$ is an inner $\N(c)$-block, then $u\in\supp a\cap K_0$ and $\beta(u)=Z_{K_0}$; if also $T_{l+1}\subseteq K_0$ then $\beta(v)=Z_{K_0}$; if not, let $T_{l'}$ ($l'>l+1$) be the next block of $t$ on the chain contained in $K_0$ (it exists, as $T_r\subseteq K_0$): the block $C_{l'-1}$ meets $T_{l'-1}\not\subseteq K_0$ and $T_{l'}\subseteq K_0$, so it meets two classes of $J_a$ and is therefore an inner $\N(c)$-block, and the point $v'\in T_{l'}\cap C_{l'-1}$ lies in $\supp a\cap K_0$, so $\beta(v')=Z_{K_0}=\beta(u)$. Since consecutive points of the chain on a common block $T_l\subseteq K_0$ have blocks $\beta$ joined by $T_l$ in $J_b$, we obtain, walking along the chain and skipping the excursions outside $K_0$, a chain in $J_b$ from $\beta(x_0)=B_1$ ($x_0\in T_1\cap C_0$) to $\beta(y_r)=B_2$ ($y_r\in T_r\cap C_r$). So $B_1,B_2$ lie in one class of $J_b$, as required.

\emph{Conclusion.} By Step~5 the blocks of $\theta_c(s)$ inside $\supp c$ are the sets $W_L$, transversal iff $L$ meets $\supp t$, and its other blocks are the $\N(c)$-blocks; so $\theta_c(s)=\theta_c(t)=w$, which is \eqref{eq:lemD}, and $ata=btb$ by Step~2.
\end{proof}

\begin{remark}\label{rem:V}
(1) The hypothesis $bcb=b$ enters the proof of \eqref{eq:lemD} only through the facts (1) and (2) of Step~4. (2) The algebraic part shows more generally: in any projection algebra satisfying \eqref{eq:Q1}, if $bcb=b$, $ata\le b$, $ctc\le a$ and $c(btb)c=ctc$, then $ata=btb$. (3) In every projection algebra, \Vcond\ is equivalent to the three-atom clause $\theta_a(t)\le\theta_d(c),\ \theta_c(t)\le a\Rightarrow\theta_a(t)\le\theta_{\theta_d(c)}(t)$ (Prover9): the half-matching hypothesis $bcb=b$ is absorbed by writing $b=\theta_d(c)$, since $\theta_{\theta_d(c)}(c)=\theta_d(c)$ is an identity of projection algebras (Lemma~\ref{lem:friends}) and every $b$ with $bcb=b$ is of this form. For $d=c$ the clause is a consequence of \eqref{eq:Q1} ($\theta_a(t)$ is then a common lower bound of $a$ and $c$, so $ac=ca$ and $\theta_a(t)=\theta_{ac}(t)=\theta_c(t)$), and the witness $\mathfrak P_8^7$ satisfies this special case while violating \Vcond.
\end{remark}

\subsection{The condition \Pone}\label{subsec:P1}
\todo{Port Theorem 5.25 (P$_1$) and its proof.}

\subsection{The condition \Qcond\ and its corollaries \Ptwo, \Wire}\label{subsec:Q}
\todo{Port Theorem 5.29 (Q) and its proof; Corollary (Q) implies (P$_2$) implies (S) (label cor:QS); the witness $\mathfrak P$ of order $6$ and the example $\mathfrak Q$ from the notes, Sections 5.1--5.3.}

\subsection{The condition \Kfour}\label{subsec:K4}
\todo{Port Theorem 5.49 (Steps 1--7, labels $(\sigma,\gamma,\lambda)$, the graph $\mathcal L$) and the derivation of \Esix\ from it (Prover9).}

\subsection{The condition \Kzero}\label{subsec:K0}

\begin{theorem}\label{thm:K0}
Let $a,b,c,y,z$ be projections of $\P_X$ with $a\mt c$, $c\mt b$, $yc=cy$, $\theta_y(b)=\theta_b(a)$, $\theta_b(z)=b$, $\theta_z(a)=\theta_z(b)$ and $\theta_y(z)=\theta_a(b)$. Then $a\mt b$. The same holds with the hypothesis $\theta_z(a)=\theta_z(b)$ replaced by $\theta_a(z)=a$.
\end{theorem}

The second statement is a different quasi-identity, which we denote \Kzero$''$: it is not a consequence of \Kzero\ and the other conditions of this paper in projection algebras (a projection algebra of order $9$ separates them, Section~\ref{sec:comp}), while conversely \Kzero$''$ implies \Kzero\ in every finite projection algebra, by finite descent (Corollary~\ref{cor:K0pp}); the block proof below gives both at once.

\begin{proof}
We use the notation of Section~\ref{subsec:PX}: $A_i,C_i,B_i$ are the $\T$-blocks of $a,c,b$ indexed through $c$, $\Gamma_i\ni A_i,C_i$ and $\Delta_i\ni C_i,B_i$ are the corresponding classes, $u=\theta_a(b)$, $v=\theta_b(a)$, and $m=\rk a$, $n=\rk u$. Every block of $c$ meeting $A_i$ lies in $\Gamma_i$, every block of $b$ meeting $C_i$ lies in $\Delta_i$, and no $\T$-block of one of $a,b,c$ is an $\N$-block of a matched partner (Lemma~\ref{lem:PX}(c)). By Lemma~\ref{lem:PX}(a), $a\mt b$ iff $n=m$ and every class of $\varepsilon_a\vee\varepsilon_b$ meeting $\supp a$ meets $\supp b$.

\smallskip\emph{Step 1: algebra, and the cells.} From $\theta_y(u)=\theta_y(\theta_y(z))=\theta_y(z)=u$ and $\theta_y(v)=\theta_y(\theta_y(b))=v$ we get $u\le y$ and $v\le y$; as $u\le a$ and $v\le b$, \eqref{eq:Q1} gives $ay=ya$ and $by=yb$, so $\theta_y(a)=ya$ and $\theta_y(b)=yb=v$. Put $p=\theta_z(a)=\theta_z(b)$. Then $b=\theta_b(z)\mt\theta_z(b)=p$ (Lemma~\ref{lem:friends}), so $\rk p=m$. The $\T$-blocks of $p=\theta_z(a)$ are the sets $L\cap\supp z$ for the classes $L$ of $\varepsilon_a\vee\varepsilon_z$ meeting both $\supp a$ and $\supp z$; there are $m$ of them, each contains a $\T(a)$-block, and there are $m$ $\T(a)$-blocks; so every class of $\varepsilon_a\vee\varepsilon_z$ meeting $\supp a$ meets $\supp z$ and contains exactly one $\T(a)$-block, which by Lemma~\ref{lem:PX}(a) says $\theta_a(z)=a$. (Under the hypotheses of \Kzero$''$ this is given, and from here on the two conditions are treated together; only $\theta_a(z)=a$ and $\theta_b(z)=b$ will be used.) Now $(ya)z(ya)=y(aza)y=yay=ya$ and $(ya)z(ya)=a(yzy)a=aua=u$, so $\theta_y(a)=ya=u$.

Let $L_i$ be the class of $\varepsilon_a\vee\varepsilon_z$ containing $A_i$ and $L'_i$ the class of $\varepsilon_b\vee\varepsilon_z$ containing $B_i$; the $L_i$ are distinct ($\theta_a(z)=a$) and so are the $L'_i$ ($\theta_b(z)=b$). A \emph{cell} is a nonempty set $L\cap L'$ with $L$ a class of $\varepsilon_a\vee\varepsilon_z$ and $L'$ a class of $\varepsilon_b\vee\varepsilon_z$; the cells partition $X$, every block of $z$ lies in a cell, every block of $a$ lies in one class of $\varepsilon_a\vee\varepsilon_z$, and every block of $b$ in one class of $\varepsilon_b\vee\varepsilon_z$. The \emph{diagonal} cells are $D_i=L_i\cap L'_i$ (possibly empty).

\smallskip\emph{Step 2: $y$ against $a,b,c$.} Lemma~\ref{lem:PX}(b) applies to the commuting pairs $(y,a)$, $(y,b)$, $(y,c)$. We claim that for $(y,a)$ and $(y,b)$ there is no class of type (iii): \emph{every $\T(a)$-block and every $\T(b)$-block meets $\supp y$.} Let $N$ be an $\N(y)$-block which is a union of $\T(a)$-blocks, $N\supseteq A_i$. As an $\N(y)$-block, $N$ is a nontransversal block of $\theta_y(b)=v=\theta_b(a)$, i.e.\ an $\N(b)$-block or a set $M\cap\supp b$ for a class $M$ of $\varepsilon_a\vee\varepsilon_b$ missing $\supp a$; the latter is impossible as $N\supseteq A_i$, so $N$ is an $\N(b)$-block. For $(y,c)$, $N$ is of type (iii) or (iv). If $N$ is an $\N(c)$-block, it meets $A_i$, so $N\subseteq\Gamma_i$, and $N\subseteq\supp a$ gives $N=A_i$, a $\T(a)$-block that is an $\N(c)$-block, impossible. So $N$ is a union of $\T(c)$-blocks; being a block of $b$ it lies in one class $\Delta_l$, which contains the single $\T(c)$-block $C_l$, so $N=C_l$ is a $\T(c)$-block that is an $\N(b)$-block, impossible. Next let $N$ be an $\N(y)$-block which is a union of $\T(b)$-blocks, $N\supseteq B_k$. It is a nontransversal block of $v=\theta_b(a)$ which is not an $\N(b)$-block, so $N=M\cap\supp b$ for a class $M$ of $\varepsilon_a\vee\varepsilon_b$ missing $\supp a$; thus $N$ contains no $\T(a)$-block, and by what was just proved $N$ is an $\N(a)$-block. For $(y,c)$, if $N$ were an $\N(c)$-block then $N\subseteq\Delta_k$ and $N=B_k$, impossible; so $N$ is a union of $\T(c)$-blocks. Every $\T(b)$- and every $\T(c)$-block lies in some $\Delta_l$, so $N=\bigsqcup_l(N\cap\Delta_l)$, where $N\cap\Delta_l$ is both a union of $\T(b)$-blocks (hence $B_l$ or empty) and a union of $\T(c)$-blocks (hence $C_l$ or empty): $N$ is a disjoint union of blocks $B_l=C_l$. Being a block of $a$, $N$ lies in one $\Gamma_l$, so $N=B_k=C_k$ is an $\N(a)$-block which is a $\T(c)$-block, impossible.

Consequences. Every $\N(y)$-block is an $\N(a)$-block and an $\N(b)$-block (type (iv) for $(y,a)$ and $(y,b)$); every $\T(y)$-block $Y$ lies in a class of type (i) or (ii), so $Y$ lies in a union of $\T(a)$-blocks or in an $\N(a)$-block, in a union of $\T(b)$-blocks or in an $\N(b)$-block, and in a union of $\T(c)$-blocks or in an $\N(c)$-block; and every $\N$-block of $a$, $b$ or $c$ meeting a $\T(y)$-block contains it (types (ii) and (iv)). The $\T$-blocks of $u=ya$ are the classes $U_1,\dots,U_n$ of type (i) of $\varepsilon_y\vee\varepsilon_a$, each a union of $\T(a)$-blocks and of $\T(y)$-blocks, and its $\N$-blocks are the $\N(a)$-blocks; comparing with the description of $u=\theta_a(b)$ (Section~\ref{subsec:PX}), the $U_j$ are the sets $M_j\cap\supp a$ for the classes $M_1,\dots,M_n$ of $\varepsilon_a\vee\varepsilon_b$ meeting $\supp a$, and every $M_j$ meets $\supp b$. Put $I_j=\{i:A_i\subseteq U_j\}$; the $I_j$ partition $[m]$, and $a\mt b$ iff $n=m$ iff every $|I_j|=1$.

\smallskip\emph{Step 3: the classes of $\varepsilon_y\vee\varepsilon_z$.} By Lemma~\ref{lem:PX}(a) the blocks of $\theta_y(z)$ inside $\supp y$ are the sets $K\cap\supp y$ for the classes $K$ of $\varepsilon_y\vee\varepsilon_z$ meeting $\supp y$, transversal iff $K$ meets $\supp z$; the blocks of $u=ya$ inside $\supp y$ are the $U_j$ (transversal) and the $\N(a)$-blocks contained in $\supp y$ (nontransversal; type (ii)). As $\theta_y(z)=u$, the classes of $\varepsilon_y\vee\varepsilon_z$ meeting $\supp y$ are $K_1,\dots,K_n$ with $K_j\cap\supp y=U_j$ and $K_j\cap\supp z\ne\emptyset$, and, for each $\N(a)$-block $F\subseteq\supp y$, a class $K_F$ with $K_F\cap\supp y=F$. In particular each $A_i$ lies in $K_j$ for $i\in I_j$, and each $\N(a)$-block is an $\N(y)$-block or lies in its $K_F$.

\smallskip\emph{Step 4: blocks of $y$ and cells.} Call a $\T(y)$-block \emph{crossing} if it meets two $\T(a)$-blocks.

(a) \emph{A crossing block $Y$ satisfies $Y\subseteq\bigcup_i(A_i\cap B_i\cap C_i)$, the union over the $i$ with $Y\cap A_i\ne\emptyset$; in particular $Y$ lies in the union of the diagonal cells.} $Y$ meets $\supp a$, so it lies in a type (i) class $U_j$ and $Y\subseteq\supp a$. $Y$ is not inside an $\N(c)$-block $F'$: such an $F'$ would meet $A_i$ and $A_{i'}$ for two $i\ne i'$, so $F'\subseteq\Gamma_i\cap\Gamma_{i'}$. So $Y\subseteq\bigcup_kC_k$; a point of $Y\cap A_i$ lies in some $C_k\subseteq\Gamma_k$ and in $A_i\subseteq\Gamma_i$, so $k=i$: $Y\cap A_i\subseteq C_i$. A point $x\in Y\cap C_i$ lies in a block of $b$ inside $\Delta_i$, namely $B_i$ or an $\N(b)$-block $G\subseteq\Delta_i$; in the second case $G\supseteq Y$, but $Y\cap A_{i'}\subseteq C_{i'}\subseteq\Delta_{i'}$ for a second index $i'$. So $Y\cap A_i\subseteq A_i\cap B_i\cap C_i$, and a point of $A_i\cap B_i$ lies in $L_i\cap L'_i=D_i$.

(b) \emph{A non-crossing $\T(y)$-block $Y$ lies in one cell.} If $Y$ meets no $\T(a)$-block then $Y$ lies in an $\N(a)$-block, and if it meets exactly one, $A_i$, then it lies in $U_j\subseteq\supp a$ and $Y\subseteq A_i$; either way $Y$ lies in one block of $a$, hence in one class of $\varepsilon_a\vee\varepsilon_z$ and in one class of $\varepsilon_a\vee\varepsilon_c$. For $\varepsilon_b\vee\varepsilon_z$: if $Y$ lies in an $\N(b)$-block we are done; otherwise $Y\subseteq\bigcup_kB_k$, and we claim that $Y$ meets only one $B_k$. Suppose $Y$ meets $B_k$ and $B_l$, $k\ne l$. $Y$ is not inside an $\N(c)$-block (it would lie in $\Delta_k\cap\Delta_l$), so $Y\subseteq\bigcup C$'s, and as in (a) $Y\cap B_k\subseteq C_k\subseteq\Gamma_k$, $Y\cap B_l\subseteq C_l\subseteq\Gamma_l$, contradicting the fact that $Y$ lies in one class of $\varepsilon_a\vee\varepsilon_c$. So $Y\subseteq B_k\subseteq L'_k$.

(c) \emph{Every $\N(y)$-block and every block of $z$ lies in one cell}: an $\N(y)$-block is a block of $a$ and of $b$ (Step~2).

(d) \emph{Every class $K$ of $\varepsilon_y\vee\varepsilon_z$ containing no crossing block lies in one cell, and every class containing a crossing block lies in $\bigcup_iD_i$.} $K$ is connected by chains of blocks of $y$ and $z$. By (b), (c) a non-crossing block lies in one cell, so it does not leave the cell of the block before it; by (a) a crossing block lies in $\bigcup D_i$, and a non-crossing block meeting it lies in some $D_i$. Starting the chains at a crossing block, resp.\ noting that without crossing blocks no chain changes cell, gives the claim.

\smallskip\emph{Step 5: $\varepsilon_a\le\varepsilon_b\vee\varepsilon_z$.} Let $A_i\subseteq U_j\subseteq K_j$. If $K_j$ contains no crossing block, $A_i$ lies in one cell. Otherwise $K_j\subseteq\bigcup_{i'}D_{i'}$, and $A_i\subseteq L_i$ with $D_{i'}\subseteq L_{i'}$ and the $L$'s distinct, so $A_i\subseteq D_i\subseteq L'_i$. An $\N(a)$-block $F\subseteq\supp y$ lies in $K_F$, which does not meet $\supp a$ and so contains no crossing block; so $F$ lies in one cell. An $\N(a)$-block not inside $\supp y$ is an $\N(y)$-block, in one cell by Step~4(c). So every block of $a$ lies in one class of $\varepsilon_b\vee\varepsilon_z$, and therefore $\varepsilon_a\vee\varepsilon_b\le\varepsilon_b\vee\varepsilon_z$.

\smallskip\emph{Step 6: the contradiction.} Suppose $n<m$, so that some $I_j$ contains two indices, say $1,2\in I_j$. The class $U_j$ of type (i) is connected by chains alternating $\T(a)$- and $\T(y)$-blocks, so some $\T(y)$-block of $U_j$ meets two $\T(a)$-blocks: $K_j$ contains a crossing block, and by Step~5, $A_1\subseteq L'_1$ and $A_2\subseteq L'_2$, two distinct classes of $\varepsilon_b\vee\varepsilon_z$. But $A_1,A_2\subseteq M_j$, a single class of $\varepsilon_a\vee\varepsilon_b$, which by Step~5 lies in one class of $\varepsilon_b\vee\varepsilon_z$. Hence $n=m$, every class of $\varepsilon_a\vee\varepsilon_b$ meeting $\supp a$ meets $\supp b$ (Step~2), and $a\mt b$.
\end{proof}

\begin{corollary}\label{cor:K0pp}
In every finite projection algebra, \Kzero$''$ implies \Kzero.
\end{corollary}

\begin{proof}
Assume the hypotheses of \Kzero\ and put $p=\theta_z(a)=\theta_z(b)$. Then $b=\theta_b(z)\mt\theta_z(b)=p$ by Lemma~\ref{lem:friends}, and $\theta_p(a)=\theta_{\theta_z(a)}(a)=\theta_z(a)=p$; so $a\mt c\mt b\mt p$ and Lemma~\ref{lem:descent}(c) gives $\theta_a(p)=a$, i.e.\ $\theta_a(\theta_z(a))=\theta_a(z)=a$. Now \Kzero$''$ applies.
\end{proof}

\todo{Here or in Section~\ref{sec:comp}: the remark that Mace4 finds no finite model of \Qone$+$\Kzero$''+\neg$\Kzero\ at any size, by the corollary, while Prover9 does not derive it; and the Prover9 derivation of \Eone\ from \Qone, (W$_2$), \Kzero.}

\subsection{The census of order $8$ and the conditions \Eone, \Esix}\label{subsec:census}
\todo{Port the definitions of $\mathfrak P$, $\mathfrak P_7$, $\mathfrak P_8^1,\dots,\mathfrak P_8^7$ (tables or presentations), Proposition 5.19 (which condition each violates), Proposition (K0 $\Rightarrow$ E$_1$, K4 $\Rightarrow$ E$_6$), and the independence proposition 5.32.}

%==============================================================================
\section{The chain conditions and the counting invariant}\label{sec:chain}
\todo{Phase 2. Port Sections 5.14 (the identification graph, the conditions (W$_n$)), 5.20 (Lemma 5.40 invariance, Corollary 5.41 separation, Lemmas 5.42--5.43 contraction and value, Theorem 5.44 end step, Theorem 5.45 (W$_\infty$)), 5.15 (the witnesses $\mathfrak P^{\mathrm{ch}}_n$ and Theorem 5.46 no finite basis), with the radius lemma 5.39 as a remark; (R$'_n$) as a remark with its witnesses and the statement that it is verified for small $m$ only.}

%==============================================================================
\section{Graph semigroups: sufficiency}\label{sec:graph}
\todo{Phase 2. Theorems 5.8 and 5.12.}

%==============================================================================
\section{Computations}\label{sec:comp}
\todo{Phase 3. The enumeration of regular $*$-semigroups to order $8$ and of projection algebras to order $9$, the embedding searches, the candidates, the SAT verifications, the Prover9/Mace4 runs (k4.pf1, k4.pf2, k0\_implies\_e1, k4\_implies\_e6, rpp\_vs\_chain, k6\_vs\_chain), the order-9 model separating \Kzero$''$ from \Kzero.}

%==============================================================================
\section{Open problems}\label{sec:open}
\todo{The three order-9 candidates and the conditions K1c/K1d, K3b that explain them; the chain family (R$'_n$); whether the list is sufficient for order $9$; the status of \Kzero$''$ versus \Kzero\ in infinite projection algebras; finite-algebraic twins of the geometric arguments.}

\bibliographystyle{amsplain}
\begin{thebibliography}{9}
\bibitem{EGMR} J.~East, R.~Gray, P.~A.~Azeef Muhammed, N.~Ru\v skuc, \emph{Projection algebras and free regular $*$-semigroups}, \todo{details}.
\bibitem{EM} J.~East, P.~A.~Azeef Muhammed, \todo{arXiv:2301.04845, title}.
\bibitem{Imaoka} T.~Imaoka, \todo{1983 paper on projection algebras / regular $*$-semigroups}.
\bibitem{Jones} P.~R.~Jones, \emph{A common framework for restriction semigroups and regular $*$-semigroups}, \todo{J. Pure Appl. Algebra 2012}.
\bibitem{EMRT} J.~East, J.~D.~Mitchell, N.~Ru\v skuc, M.~Torpey, \emph{Congruence lattices of finite diagram monoids}, Adv. Math. (2018).
\end{thebibliography}

\end{document}
